\chapter{Podstawy algebry abstrakcyjnej}
\label{ch:algebra-abstrakcyjna}
\index{algebra abstrakcyjna@Algebra abstrakcyjna}

Zbierzemy pojęcia grup, pierścieni i modułów potrzebne
do algebry homologicznej. Konstrukcje ilorazowe pozwalają
narzucać relacje, a iloczyn tensorowy pozwala zastąpić
odwzorowanie dwuliniowe liniowym. Wcześniejszy
Rozdział~\ref{ch:algebra-tensorowa} dotyczył przestrzeni
nad ciałem; tutaj pracujemy nad pierścieniem.
W Rozdziale~\ref{ch:grass-spin} wystąpił też moduł
algebry Clifforda; wrócimy do tego przykładu.

\section{Grupy, słowa i prezentacje}
\label{sec:grupy-prezentacje-5}
\index{grupa wolna@Grupa wolna}
\index{prezentacja grupy@Prezentacja grupy}

Zanim przejdziemy do modułów, zbierzmy język
potrzebny przy grupie podstawowej i
Twierdzeniu~\ref{thm:van-kampen}. Grupa opisuje
odwracalne składanie działań. W grupie klas
pętli opartych w jednym punkcie działaniem
jest składanie pętli, a odwrotność reprezentuje
pętla przebiegnięta w przeciwnym kierunku.

\subsection*{Grupy, homomorfizmy i ilorazy}

\begin{definicja}[Grupa i homomorfizm]
\index{grupa i homomorfizm@Grupa i homomorfizm}
\label{def:grupa-hom-5}
\emph{Grupa} $(G,\cdot,1)$ to zbiór z łącznym
mnożeniem, elementem neutralnym $1$ i
odwrotnością $g^{-1}$ każdego $g\in G$.
Nie zakładamy przemienności. Odwzorowanie
$\varphi:G\to H$ jest \emph{homomorfizmem},
jeśli $\varphi(gh)=\varphi(g)\varphi(h)$.
Wtedy $\varphi(1)=1$ i
$\varphi(g^{-1})=\varphi(g)^{-1}$.
Istotnie, $\varphi(1)^2=\varphi(1)$, więc skracanie daje
$\varphi(1)=1$; następnie stosujemy $\varphi$ do $gg^{-1}=1$.
Bijektywny homomorfizm nazywamy \emph{izomorfizmem}.
\end{definicja}

Przykładem jest $(\Z,+,0)$;
w zapisie addytywnym odwrotnością $m$
jest $-m$. Grupa permutacji trzech
elementów dostarcza przykładu
nieprzemiennego: przestawienia
$(12)(23)=(123)$ i $(23)(12)=(132)$ są różne
(permutacje składamy od prawej do lewej).

\begin{definicja}[Podgrupa normalna i iloraz]
\index{podgrupa normalna i iloraz@Podgrupa normalna i iloraz}
\label{def:normal-quotient-5}
\emph{Podgrupa} $N\leq G$ to podzbiór zawierający $1$
i zamknięty na mnożenie oraz branie odwrotności.
\emph{Warstwa lewostronna} to $gN=\{gn:n\in N\}$.
Mamy $gN=hN$ dokładnie wtedy, gdy $h^{-1}g\in N$:
z równości warstw wynika $g=hn$ dla pewnego $n\in N$,
a z $g=hn$ wynika $gN=hnN=hN$.
Podgrupa $N\leq G$ jest \emph{normalna},
gdy $gNg^{-1}=N$ dla każdego $g\in G$.
Wówczas zbiór warstw $G/N$ staje się
grupą przez $(gN)(hN)=ghN$.
Jądro homomorfizmu
$\ker\varphi=\{g:\varphi(g)=1\}$
jest normalne. \emph{Domknięcie normalne}
zbioru $R\subset G$, oznaczane
$\langle\!\langle R\rangle\!\rangle$,
jest najmniejszą podgrupą normalną
zawierającą $R$.
\end{definicja}

\begin{proof}[Dlaczego iloraz jest dobrze określony?]
Jeśli $g'=gn$ i $h'=hm$ dla $n,m\in N$,
to $g'h'=gnhm=gh(h^{-1}nh)m$.
Normalność daje $h^{-1}nh\in N$,
zatem $g'h'N=ghN$. Mnożenie
nie zależy więc od wybranych
reprezentantów. Łączność dziedziczymy z $G$;
elementem neutralnym jest $N$, a odwrotnością $gN$ jest $g^{-1}N$.
Dla jądra mamy $\varphi(1)=1$, a jeśli $x,y\in\ker\varphi$,
to $\varphi(xy^{-1})=\varphi(x)\varphi(y)^{-1}=1$;
jest więc podgrupą. Ponadto
$\varphi(gng^{-1})=
\varphi(g)1\varphi(g)^{-1}=1$.
Elementy domknięcia normalnego
są skończonymi iloczynami
sprzężeń $grg^{-1}$ i ich
odwrotności, $r\in R$:
zbiór tych iloczynów jest
podgrupą normalną zawierającą
$R$ i mieści się w każdej takiej
podgrupie.
\end{proof}

Normalność jest dokładnie warunkiem potrzebnym do mnożenia warstw.
Istotnie, gdyby iloczyn warstw był niezależny od reprezentantów,
to dla $n\in N$ mielibyśmy
$(gN)(nN)(g^{-1}N)=(gN)N(g^{-1}N)=N$,
czyli $gng^{-1}\in N$. Zmieniając $g$ na $g^{-1}$,
dostajemy równość $gNg^{-1}=N$.
Na przykład w $S_3$ podgrupa $N=\{1,(12)\}$ nie jest normalna:
$(23)(12)(23)=(13)\notin N$. Przy próbie obliczenia
$N\cdot(23)N$ wybór reprezentanta $1$ albo $(12)$ pierwszej
warstwy daje odpowiednio $(23)N$ i $(12)(23)N$.
Są to różne warstwy, bo $(23)^{-1}(12)(23)=(13)\notin N$.

\begin{stwierdzenie}[Twierdzenie o homomorfizmie grup]
\index{twierdzenie o homomorfizmie grup}
\label{prop:group-first-isomorphism-5}
Jeśli $\varphi:G\to H$ jest homomorfizmem i $N=\ker\varphi$,
to $\varphi$ wyznacza izomorfizm
\[
 \overline\varphi:G/N\longrightarrow\operatorname{im}\varphi,
 \qquad gN\longmapsto\varphi(g).
\]
\end{stwierdzenie}
\begin{proof}
Jądro $N$ jest normalne, więc $G/N$ jest grupą. Jeśli $gN=hN$,
to $h^{-1}g\in N$, skąd
$\varphi(h)^{-1}\varphi(g)=\varphi(h^{-1}g)=1$.
Wzór na $\overline\varphi$ nie zależy zatem od reprezentanta.
Ponadto
$\overline\varphi((gN)(hN))=\varphi(gh)
=\varphi(g)\varphi(h)$, więc jest homomorfizmem.
Jego obrazem jest z definicji $\operatorname{im}\varphi$.
Jeśli $\overline\varphi(gN)=\overline\varphi(hN)$, to
$\varphi(h^{-1}g)=1$, czyli $h^{-1}g\in N$ i $gN=hN$.
Stąd $\overline\varphi$ jest iniektywne. Równoważnie, dwa
elementy $G$ mają ten sam obraz dokładnie wtedy, gdy należą
do tej samej warstwy jądra.
\end{proof}

\begin{przyklad}[Reszty modulo $m$]
\label{prz:Zmodm-grupa-5}
Niech $m\geq2$.
W grupie addytywnej $\Z$ każda
podgrupa jest normalna.
Iloraz $\Z/m\Z$ ma elementy
$[0],\ldots,[m-1]$ i dodawanie
$[a]+[b]=[a+b]$. Homomorfizm
$\Z\to\Z/m\Z$ ma jądro $m\Z$.
Na przykład w $\Z/5\Z$ mamy $[3]+[4]=[2]$, ponieważ
$3+4-2=5\in5\Z$. Twierdzenie~\ref{prop:group-first-isomorphism-5}
wyjaśnia, że iloraz utożsamia dokładnie liczby o tej samej
reszcie modulo $m$.
\end{przyklad}

\subsection*{Grupy wolne i iloczyn wolny}

\begin{definicja}[Grupa wolna i iloczyn wolny]
\index{grupa wolna i iloczyn wolny@Grupa wolna i iloczyn wolny}
\index{slowo w grupie wolnej@słowo w grupie wolnej}
\label{def:free-groups-5}
Niech $A$ będzie zbiorem \emph{symboli},
zwanych generatorami. Dla każdego $a\in A$
tworzymy nowy symbol $a^{-1}$;
zbiór tych nowych symboli oznaczamy $A^{-1}$
i przyjmujemy $(a^{-1})^{-1}=a$.
\emph{Słowo} w alfabecie
$A^{\pm1}=A\sqcup A^{-1}$ jest skończonym
ciągiem liter $x_1\cdots x_r$, $r\geq0$.
Długość słowa to $r$; jedyne słowo długości
zero oznaczamy $\varnothing$ i nazywamy pustym.
Słowa można sklejać, czyli konkatenować:
$(x_1\cdots x_r)(y_1\cdots y_s)
=x_1\cdots x_ry_1\cdots y_s$.

\emph{Redukcja elementarna} usuwa sąsiednie
litery $xx^{-1}$ lub $x^{-1}x$.
Słowo jest \emph{zredukowane}, gdy nie zawiera
takiej pary. Dwa słowa uznajemy za równoważne,
jeżeli można przejść od jednego do drugiego
skończonym ciągiem wstawień lub usunięć
takich par. \emph{Grupa wolna} $F(A)$
to zbiór tych klas, z mnożeniem zadanym
przez sklejanie słów.

Dla grup $G,H$ \emph{iloczyn wolny} $G*H$
powstaje ze słów w \emph{rozłącznej sumie} alfabetów
$(G\setminus\{1_G\})\sqcup(H\setminus\{1_H\})$.
Litera pamięta swój czynnik, nawet gdy $G=H$.
Redukcja zastępuje sąsiednie litery jednego czynnika
ich iloczynem w tym czynniku; jeśli iloczyn jest jedynką,
usuwa obie litery. Równoważność słów dopuszcza zarówno
te redukcje, jak i operacje odwrotne.
$G*H$ jest zbiorem klas z mnożeniem przez konkatenację.
Słowo zredukowane ma litery z obu czynników naprzemiennie.
\end{definicja}

Zanim utożsamimy klasy ze słowami zredukowanymi,
trzeba wiedzieć, że wynik redukcji nie zależy
od kolejności kasowania par. Sam zbiór
\emph{wszystkich} słów z konkatenacją
jest tylko monoidem: słowa $aa^{-1}$
i $\varnothing$ nie są w nim dosłownie równe.

\begin{stwierdzenie}[Jednoznaczna postać zredukowana]
\index{redukcja slowa@redukcja słowa}
\label{prop:free-reduction-5}
Każda klasa w $F(A)$ zawiera dokładnie jedno
słowo zredukowane. Jeżeli $\operatorname{red}(w)$
oznacza to słowo, działanie grupowe ma postać
\[
 u\cdot v=\operatorname{red}(uv),\qquad
 1=\varnothing,\qquad
 (x_1\cdots x_r)^{-1}=x_r^{-1}\cdots x_1^{-1}.
\]
\end{stwierdzenie}
\begin{proof}
Czytamy słowo od lewej do prawej i trzymamy
\emph{stos} już przeczytanych liter.
Gdy następna litera jest odwrotnością
litery na wierzchu stosu, zdejmujemy wierzch;
w przeciwnym razie dokładamy nową literę.
Po ostatnim kroku na stosie nie ma sąsiedniej
pary odwrotnej, więc otrzymaliśmy słowo
zredukowane. Po każdym kroku stos reprezentuje tę samą
klasę co przeczytany fragment: dopisujemy literę albo
usuwamy końcową parę odwrotną. Wynik należy więc do klasy
początkowego słowa.

Sprawdźmy niezależność od usuwania par.
Operacja dopisania litery $x$ do słowa
na stosie jest odwracalna: operacja
dopisania $x^{-1}$ cofa ją, zarówno gdy
pierwszy krok dołożył literę, jak i wtedy,
gdy ją skasował. Dlatego przeczytanie
fragmentu $xx^{-1}$ nie zmienia stanu stosu.
Usunięcie takiego fragmentu z dowolnego miejsca
słowa nie zmienia końcowego wyniku algorytmu.
Każde słowo równoważne daje zatem ten sam wynik.
Jeśli słowo jest już zredukowane, algorytm
pozostawia je bez zmian; w każdej klasie
może więc być tylko jedno takie słowo.

Sklejanie jest łączne przed redukcją, a
usuwanie par nie zmienia wyniku algorytmu.
Stąd $\operatorname{red}(\operatorname{red}(uv)w)
=\operatorname{red}(uvw)
=\operatorname{red}(u\operatorname{red}(vw))$.
Słowo puste jest jedynką. Zapisane w tezie
słowo odwrotne po sklejeniu z oryginałem
redukuje się kolejno od środka do słowa pustego.
To sprawdza wszystkie aksjomaty grupy.
\end{proof}

\begin{przyklad}[Słowa i działanie na słowach]
\label{prz:redukcja-slow-5}
W $F(\{a,b\})$ słowo
$ab\,b^{-1}a^{-1}ba$
redukujemy najpierw do
$aa^{-1}ba$, a następnie
do $ba$. Słowa $ab$ i $ba$
są już zredukowane i różne,
więc grupa ta nie jest przemienna.
Dla jednego generatora każde słowo
redukuje się do $a^m$ z jednoznacznym
$m\in\Z$; stąd $F(\{a\})\cong\Z$.

Grupa $F(A)$ działa na zbiorze zredukowanych
słów przez mnożenie z lewej:
$u\cdot w=\operatorname{red}(uw)$.
Na przykład $a\cdot(a^{-1}b)=b$.
Jest to działanie, bo łączność udowodniono
powyżej; odwrotna operacja do działania
przez $a$ to działanie przez $a^{-1}$.
Jest ono przechodnie, gdyż $w\cdot\varnothing=w$,
i wolne: jeśli $u\cdot w=w$, to w grupie
$uw=w$, zatem po pomnożeniu przez $w^{-1}$
dostajemy $u=1$. Ten sam zbiór ma również
działanie z prawej $w\cdot u=\operatorname{red}(wu)$.
Rysunek~\ref{fig:free-words-5} pokazuje
pojedyncze kroki tego drugiego działania.
Na przykład iloczyn dwóch zredukowanych słów
$ab$ i $b^{-1}a$ trzeba jeszcze skrócić:
$\operatorname{red}(ab\,b^{-1}a)=a^2$.
Każde słowo zredukowane wyznacza drogę od $1$ bez natychmiastowego
powrotu po przebytej krawędzi. Jednoznaczność postaci zredukowanej
oznacza, że takie drogi nie tworzą cykli: pełny graf jest drzewem,
a Rysunek~\ref{fig:free-words-5} pokazuje tylko jego fragment.
Istotnie, cykl prosty można przez mnożenie z lewej przenieść tak,
aby zaczynał się w $1$. Litery kolejnych krawędzi dałyby wtedy
niepuste słowo zredukowane reprezentujące jedynkę, wbrew
Stwierdzeniu~\ref{prop:free-reduction-5}.
\end{przyklad}

\begin{figure}[htbp]
\centering
\begin{tikzpicture}[>=Latex,font=\small,
 every node/.style={circle,draw=manifoldblue!75!black,
 fill=manifoldblue!8,inner sep=2pt,minimum size=8mm}]
 \node (e) at (0,0) {$1$};
 \node (a) at (2,0) {$a$};
 \node (ai) at (-2,0) {$a^{-1}$};
 \node (b) at (0,1.5) {$b$};
 \node (bi) at (0,-1.5) {$b^{-1}$};
 \node (aa) at (3.8,.7) {$a^2$};
 \node (ab) at (3.8,-.7) {$ab$};
 \node (aiai) at (-3.8,.7) {$a^{-2}$};
 \node (aib) at (-3.8,-.7) {$a^{-1}b$};
 \node (bb) at (-.95,2.8) {$b^2$};
 \node (ba) at (.95,2.8) {$ba$};
 \node (bibi) at (-.95,-2.8) {$b^{-2}$};
 \node (bia) at (.95,-2.8) {$b^{-1}a$};
 \foreach \x/\y in {e/a,e/ai,e/b,e/bi,a/aa,a/ab,
 ai/aiai,ai/aib,b/bb,b/ba,bi/bibi,bi/bia}
   \draw[thick,manifoldblue!70!black] (\x)--(\y);
\end{tikzpicture}
\caption{Fragment grafu słów zredukowanych w
$F(\{a,b\})$. Krawędź oznacza dopisanie
z prawej jednej litery $a^{\pm1}$ albo $b^{\pm1}$
i ewentualną redukcję. Powrót po krawędzi
odpowiada dopisaniu litery odwrotnej.}
\label{fig:free-words-5}
\end{figure}

\begin{stwierdzenie}[Własność uniwersalna grupy wolnej]
\label{prop:free-universal-5}
Dla dowolnej grupy $K$ i dowolnego wyboru
elementów $h_a\in K$, $a\in A$, istnieje
dokładnie jeden homomorfizm $\Phi:F(A)\to K$
spełniający $\Phi(a)=h_a$.
\end{stwierdzenie}
\begin{proof}
Na słowach musi on mieć postać
\[
 \Phi(a_1^{\epsilon_1}\cdots a_r^{\epsilon_r})
 =h_{a_1}^{\epsilon_1}\cdots h_{a_r}^{\epsilon_r},
 \qquad \epsilon_i\in\{1,-1\}.
\]
Usunięcie sąsiedniej pary odwrotnej nie zmienia
iloczynu po prawej. Wzór zależy więc tylko
od klasy słowa i jest dobrze określony.
Sklejanie słów odpowiada mnożeniu obrazów,
zatem $\Phi$ jest homomorfizmem.
Każde słowo jest iloczynem generatorów
i ich odwrotności, co dowodzi jednoznaczności.
Właśnie brak dodatkowych relacji między
generatorami nazywamy \emph{wolnością}.
\end{proof}

\begin{stwierdzenie}[Postać normalna i własność uniwersalna iloczynu wolnego]
\label{prop:free-product-normal-5}
Każdy element $G*H$ ma dokładnie jeden zapis jako słowo
naprzemienne o literach różnych od jedynki, z dopuszczeniem
słowa pustego. Odwzorowania $G\to G*H$ i $H\to G*H$,
wysyłające element niejednostkowy na słowo jednoliterowe,
a jedynkę na słowo puste, są iniektywnymi homomorfizmami.
Dla każdej pary homomorfizmów $f:G\to K$, $h:H\to K$
istnieje dokładnie jeden homomorfizm $\Phi:G*H\to K$
ograniczający się na czynnikach do $f$ oraz $h$.
\end{stwierdzenie}
\begin{proof}
Czytamy litery po kolei i przechowujemy stos liter
niejednostkowych, pochodzących naprzemiennie z obu czynników.
Dla $x$ z jednego czynnika oznaczmy przez $T_x$ operację
dopisania $x$: literę z innego czynnika dokładamy, a z tego
samego mnożymy z ostatnią literą stosu. Gdy iloczyn jest
jedynką, zdejmujemy tę ostatnią literę. Przyjmujemy $T_1=\id$.
Każdy krok jest dozwoloną redukcją, więc wynik algorytmu
należy do klasy czytanego słowa.

Kluczowa jest równość $T_y(T_x(u))=T_{xy}(u)$, gdy $x,y$
pochodzą z tego samego czynnika. Jeśli któraś litera jest
jedynką, równość jest natychmiastowa. W pozostałym przypadku:
\begin{itemize}[leftmargin=*]
\item Gdy $u$ jest pusty albo kończy się literą drugiego
czynnika, najpierw dopisujemy $x$, potem zastępujemy $xy$
jedną literą albo usuwamy je, jeśli $xy=1$. To właśnie $T_{xy}(u)$.
\item Gdy $u=vz$ i $z$ jest z tego samego czynnika co $x,y$,
słowo $v$ jest puste albo kończy się w drugim czynniku.
Jeśli $zx\ne1$, dwa kroki zastępują końcówkę przez
$(zx)y=z(xy)$, z usunięciem jedynki. Dla $xy=1$ odzyskujemy $vz$;
dla $xy\ne1$ jest to pojedynczy krok $T_{xy}$.
Jeśli $zx=1$, pierwszy krok usuwa $z$, a drugi daje $vy$.
Z kolei $z(xy)=y$; gdy $xy\ne1$, pojedynczy krok również daje $vy$.
Gdy $xy=1$, mamy $y=z$, więc $vy=u=T_1(u)$.
\end{itemize}
Zastąpienie sąsiednich liter ich iloczynem nie zmienia zatem
stosu po przeczytaniu tego fragmentu, ani wyniku po przeczytaniu
reszty słowa. Dotyczy to także odwrotnej operacji.
Wynik jest stały na klasie równoważności. Ponieważ algorytm
pozostawia słowo naprzemienne bez zmian, w każdej klasie jest
dokładnie jedno takie słowo.

Relacje można wykonywać wewnątrz dowolnego dłuższego słowa,
więc konkatenacja jest dobrze określona na klasach i łączna.
Słowo puste jest jedynką, a odwrócenie kolejności liter
i zastąpienie ich odwrotnościami daje odwrotność słowa.
Relacje wewnątrz czynników zapewniają, że ich odwzorowania
do $G*H$ są homomorfizmami. Jednoznaczność postaci normalnej
daje ich iniektywność; obrazy czynników mają tylko jedynkę wspólną.

Definiujemy $\Phi(x_1\cdots x_r)$ jako iloczyn obrazów liter,
używając $f$ lub $h$ zależnie od czynnika.
Redukcja w jednym czynniku nie zmienia tego iloczynu, bo
$f$ i $h$ są homomorfizmami. Wzór przechodzi więc na klasy.
Konkatenacja słów odpowiada mnożeniu obrazów, a każdy element
jest iloczynem liter z czynników; daje to istnienie i jednoznaczność $\Phi$.
\end{proof}

Ta własność uniwersalna będzie używana w
Twierdzeniu~\ref{thm:van-kampen}.

Na przykład dla $G=\langle a\rangle\cong\Z$
i $H=\langle b\rangle\cong\Z$ własności
uniwersalne dają $G*H\cong F(\{a,b\})$.
Iloczyn słów naprzemiennych
$(ab)(b^{-1}a)$ redukuje się do $a^2$:
po skasowaniu $bb^{-1}$ zostaje $aa$.
Podobnie, dla $G=\{1,s\}$ z $s^2=1$ i
$H=\{1,t,t^2\}$ z $t^3=1$ mamy w $G*H$ redukcję
\[
 (st^2)(ts)=s(t^2t)s=s\,1\,s=s^2=1.
\]
Pierwsze skrócenie odbywa się w $H$, a ostatnie w $G$;
oba są konieczne, by otrzymać postać normalną.

Nie należy mylić iloczynu wolnego z iloczynem prostym.
Dla $G=\langle a\mid a^2=1\rangle$ i
$H=\langle b\mid b^2=1\rangle$, obu rzędu dwa,
słowa $(ab)^k$, $k\geq0$, są różnymi postaciami normalnymi.
Zatem $G*H$ jest nieskończony, podczas gdy $G\times H$ ma
cztery elementy. W $G*H$ nie narzucamy przemienności liter
z różnych czynników.

\subsection*{Generatory i relacje}

\begin{definicja}[Prezentacja grupy]
\index{prezentacja grupy@Prezentacja grupy}
\label{def:prezentacja-grupy-5}
Jeśli $R$ jest zbiorem słów
w $F(A)$, to
\[
 \langle A\mid R\rangle
 :=F(A)/\langle\!\langle R\rangle\!\rangle
\]
nazywamy \emph{grupą o generatorach}
$A$ i \emph{relacjach} $R$.
Relacja $r=1$ oznacza, że jej
wszystkie konsekwencje po
mnożeniu i sprzęganiu również
stają się jedynką.
\end{definicja}

\begin{stwierdzenie}[Odwzorowanie zadane na generatorach i relacjach]
\label{prop:presentation-universal-5}
Niech $K$ będzie grupą. Wybór elementów $h_a\in K$ dla $a\in A$
wyznacza homomorfizm $\langle A\mid R\rangle\to K$ wysyłający
klasę $a$ na $h_a$ wtedy i tylko wtedy, gdy każde słowo $r\in R$
po podstawieniu $a\mapsto h_a$ daje jedynkę w $K$.
Jeśli taki homomorfizm istnieje, jest jedyny.
\end{stwierdzenie}
\begin{proof}
Z własności uniwersalnej grupy wolnej
(Stwierdzenie~\ref{prop:free-universal-5}) otrzymujemy jedyny
homomorfizm $f:F(A)\to K$ o zadanych obrazach generatorów.
Jeśli $f(r)=1$ dla wszystkich $r\in R$, normalność jądra
implikuje $\langle\!\langle R\rangle\!\rangle\subseteq\ker f$.
Możemy więc określić
$\bar f(w\langle\!\langle R\rangle\!\rangle)=f(w)$.
Gdy dwie warstwy są równe, ich reprezentanci różnią się elementem
jądra $f$, zatem wzór jest dobrze określony; zachowuje też iloczyn.
Odwrotnie, homomorfizm z ilorazu wysyła każdą klasę $r\in R$
na jedynkę, więc podstawione słowo musi być jedynką.
Jednoznaczność wynika z tego, że klasy liter $a$ generują iloraz.
\end{proof}

\begin{przyklad}[Jedna relacja potęgowa]
\label{prz:cyclic-presentation-5}
Dla $m\geq2$ grupa $\langle a\mid a^m\rangle$ jest izomorficzna
z $\Z/m\Z$. Każde jej słowo jest potęgą $a^k$;
relacja $a^m=1$ pozwala zmniejszyć wykładnik modulo $m$.
Nie powoduje ona dalszych utożsamień: przypisanie
$a\mapsto[1]\in\Z/m\Z$ respektuje $a^m=1$ i daje homomorfizm
z prezentacji na $\Z/m\Z$. Obrazy $1,a,\ldots,a^{m-1}$ są różne,
więc te $m$ potęg są różnymi elementami grupy.
\end{przyklad}

\begin{przyklad}[Prezentacja grupy torusa]
\label{prz:prezentacja-torusa-5}
Niech $F(\{a,b\})$ będzie grupą wolną
na dwóch symbolach. Relacja
$aba^{-1}b^{-1}=1$
jest równoważna $ab=ba$.
Samo słowo $aba^{-1}b^{-1}$ jest zredukowane i niepuste
w $F(\{a,b\})$; staje się jedynką dopiero po przejściu do ilorazu.
Po jej narzuceniu każde słowo
można przestawić do postaci
$a^m b^n$, $m,n\in\Z$.
Odwzorowanie
$F(\{a,b\})\to\Z^2$ dane przez
$a\mapsto(1,0)$,
$b\mapsto(0,1)$
zabija tę relację i wysyła
$a^m b^n$ na $(m,n)$.
Postacie te są więc wzajemnie
różne, a zatem
\[
 \langle a,b\mid aba^{-1}b^{-1}\rangle
 \cong\Z^2.
\]
Sklejenie przeciwnych boków kwadratu
tworzy torus. Brzeg kwadratu przed
sklejeniem odczytujemy przeciwnie do ruchu wskazówek zegara
jako słowo $aba^{-1}b^{-1}$; orientacje boków
pokazuje Rysunek~\ref{fig:torus-square-5}.
Topologiczne uzasadnienie, że
prezentacja opisuje $\pi_1(T^2)$,
podamy przy
Twierdzeniu~\ref{thm:van-kampen}.
\end{przyklad}

\begin{figure}[htbp]
\centering
\begin{tikzpicture}[>=Latex,font=\small]
 \fill[manifoldblue!5] (0,0) rectangle (3.2,3.2);
 \draw[very thick,manifoldblue!75!black] (0,0) rectangle (3.2,3.2);
 \draw[->,very thick,manifoldblue!75!black] (.5,0)--(2.7,0);
 \draw[->,very thick,manifoldblue!75!black] (.5,3.2)--(2.7,3.2);
 \draw[->,very thick,manifoldblue!75!black] (0,.5)--(0,2.7);
 \draw[->,very thick,manifoldblue!75!black] (3.2,.5)--(3.2,2.7);
 \node[below] at (1.6,0) {$a$};
 \node[above] at (1.6,3.2) {$a$};
 \node[left] at (0,1.6) {$b$};
 \node[right] at (3.2,1.6) {$b$};
 \foreach \x/\y in {0/0,3.2/0,3.2/3.2,0/3.2}
   \fill[manifoldblue!75!black] (\x,\y) circle (1.5pt);
 \node at (1.6,1.6) {$[0,1]^2$};
\end{tikzpicture}
\caption{Kwadrat fundamentalny torusa. Strzałki wskazują
orientacje identyfikowanych par boków. Obchodząc brzeg kwadratu
od lewego dolnego wierzchołka przeciwnie do ruchu wskazówek zegara,
czytamy $a,b,a^{-1},b^{-1}$; wszystkie cztery wierzchołki
po sklejeniu reprezentują ten sam punkt.}
\label{fig:torus-square-5}
\end{figure}


\subsection*{Iloczyn prosty grup}
\index{iloczyn prosty grup}
\index{grupa!iloczyn prosty}

Iloczyn prosty opisuje niezależne wykonywanie dwóch działań.
\begin{definicja}[Iloczyn prosty]
Dla grup $G,H$ ich \emph{iloczyn prosty} jest zbiorem $G\times H$
z mnożeniem $(g_1,h_1)(g_2,h_2)=(g_1g_2,h_1h_2)$.
Jedynką jest $(e_G,e_H)$, a $(g,h)^{-1}=(g^{-1},h^{-1})$.
Łączność wynika z łączności w każdej współrzędnej.
\end{definicja}

Inkluzje $g\mapsto(g,e_H)$ i $h\mapsto(e_G,h)$ są homomorfizmami.
Ich obrazy są normalne: sprzęganie odbywa się współrzędnie.
Elementy z różnych czynników komutują i każdy element ma jednoznaczny
rozkład $(g,h)=(g,e_H)(e_G,h)$.

\begin{stwierdzenie}[Wewnętrzny iloczyn prosty]
Jeśli $G,H$ są normalnymi podgrupami $K$, $G\cap H=\{e\}$ i $K=GH$,
to $F:G\times H\to K$, $F(g,h)=gh$, jest izomorfizmem.
Zapis $GH$ oznacza zbiór wszystkich iloczynów $gh$.
\end{stwierdzenie}
\begin{proof}
Komutator $ghg^{-1}h^{-1}$ należy do $G$ i do $H$ przez normalność:
jest iloczynem $g(hg^{-1}h^{-1})$ elementów $G$ oraz
iloczynem $(ghg^{-1})h^{-1}$ elementów $H$,
więc jest jedynką. Dlatego $gh=hg$ i
\[
 F(g_1g_2,h_1h_2)=g_1g_2h_1h_2=g_1h_1g_2h_2
 =F(g_1,h_1)F(g_2,h_2).
\]
Warunek $K=GH$ daje surjektywność. Jeśli $F(g,h)=e$, to
$g=h^{-1}\in G\cap H$, zatem $g=h=e$; jądro jest trywialne.
\end{proof}

\begin{przyklad}[Grupy cykliczne i grupa Kleina]
Grupę generowaną przez jeden element nazywamy \emph{cykliczną}.
Piszemy $C_m=\Z/m\Z$ dla $m\geq1$. Rząd elementu jest najmniejszą
dodatnią liczbą potęg dającą jedynkę, o ile taka liczba istnieje.
W $C_m\times C_n$ element $([1],[1])$ ma rząd
$\operatorname{lcm}(m,n)$: jego $j$-krotność znika dokładnie wtedy,
gdy $m\mid j$ i $n\mid j$. Dla $\gcd(m,n)=1$ rząd wynosi $mn$,
więc ten element generuje całą grupę i $C_m\times C_n\cong C_{mn}$.
Na przykład $C_2\times C_3\cong C_6$. Natomiast w
$V_4:=C_2\times C_2$, zwanej \emph{grupą Kleina}, każdy niezerowy
element ma rząd dwa, więc nie ma generatora całej czteroelementowej grupy.
\end{przyklad}

\subsection*{Iloczyn półprosty grup}
\index{iloczyn półprosty grup}
\index{grupa!iloczyn półprosty}
\index{działanie grupy!przez automorfizmy}

Teraz jedna grupa może zmieniać sposób działania drugiej.
\emph{Automorfizmem} grupy $N$ nazywamy izomorfizm $N\to N$;
ich grupa ze składaniem to $\operatorname{Aut}(N)$.
Działanie $H$ na $N$ przez automorfizmy jest homomorfizmem
$\varphi:H\to\operatorname{Aut}(N)$, a więc
$\varphi(h_1h_2)=\varphi(h_1)\circ\varphi(h_2)$.

\begin{definicja}[Iloczyn półprosty]
\emph{Iloczyn półprosty} $N\rtimes_\varphi H$ ma zbiór elementów
$N\times H$ i mnożenie
\[
 (n_1,h_1)(n_2,h_2)
 =\bigl(n_1\varphi(h_1)(n_2),h_1h_2\bigr).
\]
Gdy działanie jest ustalone, skracamy zapis do $N\rtimes H$.
\end{definicja}

Sprawdźmy łączność. Pierwsza współrzędna obu sposobów mnożenia
trzech elementów jest równa
\[
 n_1\varphi(h_1)(n_2)\varphi(h_1h_2)(n_3)
 =n_1\varphi(h_1)\bigl(n_2\varphi(h_2)(n_3)\bigr).
\]
Druga to $h_1h_2h_3$. Jedynką jest $(e_N,e_H)$, a odwrotność
ma postać $(n,h)^{-1}=\bigl(\varphi(h^{-1})(n^{-1}),h^{-1}\bigr)$.
Wyznaczmy ją z równania $(n,h)(a,k)=(e_N,e_H)$:
druga współrzędna daje $k=h^{-1}$, a pierwsza
$\varphi(h)(a)=n^{-1}$, czyli $a=\varphi(h^{-1})(n^{-1})$.
Ten wybór działa z obu stron. Pierwsze współrzędne obu iloczynów to
\[
 \begin{aligned}
 n\,\varphi(h)\bigl(\varphi(h^{-1})(n^{-1})\bigr)
 &=nn^{-1}=e_N,\\
 \varphi(h^{-1})(n^{-1})\,\varphi(h^{-1})(n)
 &=\varphi(h^{-1})(n^{-1}n)=e_N.
 \end{aligned}
\]
Drugie współrzędne wynoszą odpowiednio $hh^{-1}=e_H$ i $h^{-1}h=e_H$.
Inkluzje $n\mapsto(n,e_H)$ i $h\mapsto(e_N,h)$ są homomorfizmami.
Pierwszy obraz jest normalny jako jądro homomorfizmu
$(n,h)\mapsto h$. Drugi nie musi być normalny, a wzór
\[
 (e_N,h)(n,e_H)(e_N,h)^{-1}=(\varphi(h)(n),e_H)
\]
pokazuje, że zadane działanie realizuje się przez sprzężenie.
Dla działania trywialnego, $\varphi(h)=\id_N$, odzyskujemy iloczyn prosty.

\begin{stwierdzenie}[Wewnętrzny iloczyn półprosty]
\index{iloczyn półprosty grup!wewnętrzny}
Niech $N,H\leq G$, $N\trianglelefteq G$, $N\cap H=\{e\}$ i $G=NH$.
Wtedy $G\cong N\rtimes H$, z działaniem $\varphi(h)(n)=hnh^{-1}$.
\end{stwierdzenie}
\begin{proof}
Normalność zapewnia, że sprzężenie przez $h$ jest automorfizmem $N$,
a składanie sprzężeń daje homomorfizm $\varphi$.
Dla $F(n,h)=nh$ mamy
\[
 F\bigl(n_1(h_1n_2h_1^{-1}),h_1h_2\bigr)
 =n_1h_1n_2h_2=F(n_1,h_1)F(n_2,h_2).
\]
Surjektywność wynika z $G=NH$, a $nh=e$ wymusza
$n=h^{-1}\in N\cap H=\{e\}$, więc $F$ jest iniektywny.
\end{proof}

\begin{uwaga}[Rozszczepiony ciąg dokładny]
\index{ciąg dokładny!rozszczepiony}
\index{rozszerzenie grup!rozszczepione}
Zapis $1\to N\xrightarrow{i}G\xrightarrow{p}H\to1$ oznacza,
że $i$ jest iniekcją, $p$ surjekcją i $i(N)=\ker p$; to
\emph{krótki ciąg dokładny grup}, czyli rozszerzenie $H$ przez $N$.
Jest \emph{rozszczepiony}, gdy istnieje homomorfizm $s:H\to G$
z $p\circ s=\id_H$. Dla iloczynu półprostego są to
$i(n)=(n,e_H)$, $p(n,h)=h$ i $s(h)=(e_N,h)$.
Odwrotnie, w rozszczepionym ciągu $s$ jest iniektywny,
$i(N)\cap s(H)=\{e\}$ i każdy $g\in G$ ma rozkład
$g=(g\,s(p(g))^{-1})s(p(g))$, z pierwszym czynnikiem w $\ker p$.
Poprzednie stwierdzenie daje zatem $G\cong N\rtimes H$.
\end{uwaga}

\begin{przyklad}[Grupa diedralna]
\index{grupa diedralna!jako iloczyn półprosty}
Dla $n\geq3$ niech $r$ będzie obrotem foremnego $n$-kąta o $2\pi/n$,
a $s$ jego odbiciem. Mamy $r^n=s^2=e$ i $srs^{-1}=r^{-1}$.
Ostatnią relację można zobaczyć we współrzędnych zespolonych:
wybieramy środek wielokąta w $0$ i oś odbicia jako oś rzeczywistą.
Wtedy $r(z)=\zeta z$ dla $\zeta=e^{2\pi i/n}$,
$s(z)=\overline z=s^{-1}(z)$ i
$(s\circ r\circ s^{-1})(z)=\overline{\zeta\overline z}
=\overline\zeta z=\zeta^{-1}z=r^{-1}(z)$.
Każda symetria ma postać $r^j$ lub $r^js$, $0\leq j<n$:
jej obraz jednego wierzchołka i wybór kierunku do sąsiedniego
wyznaczają całą symetrię. Te $2n$ możliwości są różne.
Grupa symetrii $D_n$ jest więc $C_n\rtimes C_2$, gdzie nietrywialny
element $C_2$ odwraca generator $C_n$. Dla trójkąta wszystkie
permutacje trzech wierzchołków są symetriami, stąd
$S_3\cong C_3\rtimes C_2$.
Na przykład w $D_5$ wzór $sr^js^{-1}=r^{-j}$ daje
$(r^2s)(r^3s)=r^2(sr^3s)=r^{-1}$; iloczyn dwóch odbić jest obrotem.
\end{przyklad}

\begin{przyklad}[Grupa afiniczna]
\index{grupa afiniczna!jako iloczyn półprosty}
Grupa odwracalnych macierzy rzeczywistych $\mathrm{GL}(n,\R)$
działa na addytywnej grupie $\R^n$ przez $A\cdot v=Av$.
Grupa odwracalnych odwzorowań afinicznych $x\mapsto Ax+v$ to
$\operatorname{Aff}(n)=\R^n\rtimes\mathrm{GL}(n,\R)$, ponieważ
\[
 A(Bx+w)+v=ABx+(v+Aw),\qquad
 (v,A)(w,B)=(v+Aw,AB).
\]
Macierz i przesunięcie są jednoznaczne: $v$ jest obrazem zera,
a $A$ jest częścią liniową. Czynnik $Aw$ odróżnia ten iloczyn
od iloczynu prostego.
\end{przyklad}

\begin{przyklad}[Ruchy euklidesowe]
\index{grupa euklidesowa!jako iloczyn półprosty}
\index{grupa ruchów euklidesowych}
Niech $\mathrm O(n)=\{A:A^TA=I\}$ i
$\mathrm{SO}(n)=\{A\in\mathrm O(n):\det A=1\}$.
Grupa izometrii przestrzeni euklidesowej to
$E(n)=\R^n\rtimes\mathrm O(n)$, z tym samym prawem składania.
Istotnie, dla izometrii $f$ połóżmy $v=f(0)$ i $g(x)=f(x)-v$.
Mamy $g(0)=0$ i $\|g(x)-g(y)\|=\|x-y\|$. Zatem
$\|g(x)\|=\|x\|$, a tożsamość polaryzacyjna
$2\langle u,w\rangle=\|u\|^2+\|w\|^2-\|u-w\|^2$ daje
$\langle g(x),g(y)\rangle=\langle x,y\rangle$.
Wektory $g(e_i)$ tworzą bazę ortonormalną, a współrzędne $g(x)$
w tej bazie wynoszą $\langle g(x),g(e_i)\rangle=x_i$.
Zatem $g(x)=\sum_i x_i g(e_i)=Ax$ z $A\in\mathrm O(n)$.
Odwrotnie, każde $Ax+v$ z taką macierzą zachowuje odległości.
Izometrie zachowujące orientację tworzą
$E^+(n)=\R^n\rtimes\mathrm{SO}(n)$.
Więcej o grupach macierzowych znajduje się w
Dodatku~\ref{app:grupy-algebry-macierzowe}.
\end{przyklad}

\begin{uwaga}[Kolejność czynników i wybór działania]
W zapisie $N\rtimes H$ (również $H\ltimes N$) normalny jest czynnik
$N$, a $H$ działa na nim przez automorfizmy. Same grupy $N,H$ nie
wyznaczają wyniku: dla $C_3,C_2$ działanie trywialne daje abelową
grupę $C_6$, a odwracanie generatora daje nieabelową $S_3$.
Dlatego przy nieustalonym działaniu zachowujemy zapis $N\rtimes_\varphi H$.
\end{uwaga}


\section{Pierścienie, ideały i ilorazy}
\label{sec:pierscienie-idealy}
\index{pierscien@Pierścień}
\index{ideal@Ideał}

W grupie normalność podgrupy pozwala utworzyć
iloraz. Dla pierścieni podobną rolę pełni ideał:
trzeba zachować nie tylko dodawanie, lecz także
mnożenie klas przez dowolne elementy pierścienia.
W praktyce ideał jest zbiorem wyrażeń, które chcemy
uznać za zero. W ilorazie dwa elementy są równe
właśnie wtedy, gdy ich różnica należy do tego zbioru.

\begin{definicja}[Pierścień przemienny i homomorfizm]
\label{def:pierscien-przemienny}
\emph{Pierścień przemienny z jedynką} $R$ ma
dwa działania $+$ i mnożenie. Grupa $(R,+,0)$
jest abelowa, mnożenie jest łączne i przemienne,
ma jedynkę $1$, a rozdzielność
$r(s+t)=rs+rt$ wiąże oba działania.
Pierścień zerowy z $1=0$ jest dozwolony.
\emph{Homomorfizm pierścieni}
$\varphi:R\to S$ zachowuje dodawanie, mnożenie
i jedynkę: $\varphi(1_R)=1_S$.
\end{definicja}

Z przemienności mnożenia wynika też druga rozdzielność
$(r+s)t=rt+st$. Element $u\in R$ nazywamy
\emph{odwracalnym}, gdy istnieje $v\in R$ z $uv=1$;
zbiór takich elementów oznaczamy $R^\times$.
W każdym pierścieniu $0r=(0+0)r=0r+0r$,
zatem po dodaniu elementu przeciwnego $0r=0$.
Ponieważ $(r+(-r))s=0$, mamy także $(-r)s=-(rs)$.

Przykładami są $\Z$, $\Q$ i pierścień
wielomianów $\mathbb K[t]$ nad ciałem $\mathbb K$.
W $\Z$ nie każdy niezerowy element ma odwrotność
mnożeniową: równanie $2x=1$ nie ma rozwiązania
w $\Z$. Definicja pierścienia, w odróżnieniu od definicji
ciała, nie wymaga takich odwrotności. Dokładnie
$\Z^\times=\{-1,1\}$, choć $\Z_*$ zawiera wszystkie
niezerowe liczby całkowite.

\begin{definicja}[Dziedzina całkowitości i ciało]
\label{def:dziedzina-cialo-5}
\index{dziedzina całkowitości}
\index{ciało}
Niezerowy element $a$ pierścienia przemiennego jest \emph{dzielnikiem
zera}, jeśli $ab=0$ dla pewnego $b\ne0$. \emph{Dziedzina całkowitości}
to pierścień przemienny z $1\ne0$ bez dzielników zera, czyli
$ab=0$ wymusza $a=0$ lub $b=0$. Pozwala to skracać niezerowy czynnik:
z $ac=bc$ i $c\ne0$ wynika $(a-b)c=0$, więc $a=b$.
\emph{Ciało} jest pierścieniem przemiennym z $1\ne0$, w którym
każdy niezerowy element ma odwrotność mnożeniową.
\end{definicja}

Każde ciało jest dziedziną: z $ab=0$, $a\ne0$, po pomnożeniu
przez $a^{-1}$ dostajemy $b=0$. Dziedzina $\Z$ nie jest ciałem,
natomiast $\Q$ jest ciałem. Również $\mathbb K[t]$ nad ciałem
jest dziedziną, bo iloczyn niezerowych współczynników wiodących
jest niezerowy. Pierścień $\Z/6\Z$ nie jest dziedziną:
$[2][3]=[0]$, choć oba czynniki są niezerowe.
Nie wolno więc w nim skracać dowolnego niezerowego czynnika:
$[2][3]=[2][0]$, lecz $[3]\ne[0]$.

\begin{definicja}[Ideał i pierścień ilorazowy]
\label{def:ideal-iloraz}
\emph{Ideał} $I\subseteq R$ jest podgrupą
addytywną zamkniętą na mnożenie przez
elementy $R$: jeśli $r\in R$ i $x\in I$,
to $rx\in I$. Pierścień $R/I$ ma klasy
$[r]=r+I$ oraz działania
$[r]+[s]=[r+s]$, $[r][s]=[rs]$.
Klasy spełniają $[r]=[s]$ dokładnie wtedy, gdy
$r-s\in I$. Dla $a\in R$ zbiór
$(a):=aR=\{ar:r\in R\}$ jest \emph{ideałem głównym}
generowanym przez $a$.
\end{definicja}

\begin{proof}[Dlaczego iloraz jest dobrze określony?]
Jeżeli $r'=r+i$ i $s'=s+j$ dla $i,j\in I$,
to $r'+s'-(r+s)=i+j\in I$ oraz
$r's'-rs=rj+is+ij\in I$.
Dlatego wynik obu działań nie zależy od
wyboru przedstawicieli. Prawa łączności, przemienności
i rozdzielności dziedziczą się z $R$;
zerem jest $[0]$, jedynką $[1]$, a elementem
przeciwnym do $[r]$ jest $[-r]$.
Odwzorowanie ilorazowe $q:R\to R/I$, $q(r)=[r]$,
jest więc surjektywnym homomorfizmem o jądrze $I$.

Jądro dowolnego homomorfizmu $\varphi:R\to S$ jest ideałem:
z $\varphi(x)=\varphi(y)=0$ wynikają
$\varphi(x-y)=0$ i
$\varphi(rx)=\varphi(r)\varphi(x)=0$.
Gdy $I=R$, iloraz jest pierścieniem zerowym
z $[1]=[0]$.
\end{proof}

\begin{stwierdzenie}[Ideały pierścienia liczb całkowitych]
\label{prop:idealy-Z-5}
Każdy ideał $I\subseteq\Z$ ma postać $m\Z$ dla dokładnie
jednego $m\geq0$. Dla $m=0$ jest to ideał zerowy.
\end{stwierdzenie}
\begin{proof}
Jeżeli $I\ne\{0\}$, to razem z niezerowym $a$ zawiera $-a$,
więc zawiera liczbę dodatnią. Niech $m$ będzie najmniejszą taką
liczbą. Ponieważ $I$ jest zamknięty na mnożenie przez liczby
całkowite, $m\Z\subseteq I$. Dla dowolnego $a\in I$
dzielenie z resztą daje $a=qm+r$ z $0\leq r<m$.
Liczba $r=a-qm$ należy do $I$, więc z minimalności $m$
wynika $r=0$. Stąd $a\in m\Z$ i $I=m\Z$.
Najmniejszy dodatni element wyznacza $m$ jednoznacznie;
przypadek $I=\{0\}$ daje $m=0$.
\end{proof}

\begin{stwierdzenie}[Własność ilorazu pierścienia]
\label{prop:ring-quotient-universal-5}
Niech $I\subseteq R$ będzie ideałem i niech
$f:R\to S$ będzie homomorfizmem pierścieni.
Istnieje homomorfizm $\bar f:R/I\to S$ spełniający
$f=\bar f\circ q$ dokładnie wtedy, gdy $I\subseteq\ker f$.
Jeśli istnieje, jest jedyny i dany wzorem $\bar f([r])=f(r)$.
W szczególności
\[
 R/\ker f\cong\operatorname{im}f,
 \qquad [r]\longmapsto f(r).
\]
\end{stwierdzenie}
\begin{proof}
Jeżeli $f=\bar f\circ q$, to dla $i\in I$ mamy
$q(i)=[0]$, więc $f(i)=\bar f([0])=0$.
Odwrotnie, załóżmy $I\subseteq\ker f$. Z równości
$[r]=[s]$ wynika $r-s\in I$, a stąd
$f(r)-f(s)=f(r-s)=0$. Wzór na $\bar f$ jest więc
dobrze określony. Zachowuje dodawanie, mnożenie i jedynkę,
ponieważ robi to $f$; warunek $f=\bar f\circ q$ zachodzi
dla każdego $r\in R$. Surjektywność $q$ wymusza
jednoznaczność $\bar f$. Przy $I=\ker f$ obraz $\bar f$
jest równy $\operatorname{im}f$, a jego jądro jest zerowe:
$\bar f([r])=0$ oznacza $r\in\ker f$, czyli $[r]=[0]$.
Ograniczenie przeciwdziedziny do obrazu daje izomorfizm.
\end{proof}

\begin{przyklad}[Relacje w pierścieniach]
\label{prz:idealy-ilorazy}
Niech $m\geq2$. Ideał $m\Z$ narzuca
relację $m=0$, a iloraz $\Z/m\Z$ jest pierścieniem
reszt modulo $m$. Na przykład w $\Z/5\Z$ mamy
$[17]=[2]$ i $[17][4]=[68]=[3]$.
Rysunek~\ref{fig:reszty-modulo-3-5} przedstawia
klasy dla $m=3$.

W $\mathbb K[t]$
ideał $(t):=t\mathbb K[t]=
\{tf(t):f\in\mathbb K[t]\}$ składa się z wielomianów
o zerowym wyrazie stałym. Ewaluacja
$f(t)\mapsto f(0)$ ma jądro $(t)$,
a Stwierdzenie~\ref{prop:ring-quotient-universal-5} daje izomorfizm
$\mathbb K[t]/(t)\cong\mathbb K$.
Istotnie, każda klasa zawiera dokładnie
jedną stałą, bo $f(t)-f(0)$ jest podzielny
przez $t$, a jedyną stałą podzielną przez $t$ jest zero.

Inny wynik daje relacja $t^2=0$.
Każda klasa w $\mathbb K[t]/(t^2)$ ma dokładnie
jeden reprezentant $a+bt$, $a,b\in\mathbb K$,
ponieważ odrzucamy wyrazy stopnia co najmniej dwa,
a żaden niezerowy wielomian stopnia co najwyżej jeden
nie jest wielokrotnością $t^2$.
Mnożenie ma postać
$[a+bt][c+dt]=[ac+(ad+bc)t]$, ponieważ
$(a+bt)(c+dt)=ac+adt+bct+bdt^2$,
a składnik $bdt^2$ należy do ideału $(t^2)$.
Klasa $[t]$ jest niezerowa, bo $t\notin(t^2)$,
ale $[t]^2=[0]$.
Iloraz pierścienia może więc mieć dzielniki zera,
choć pierścień wyjściowy $\mathbb K[t]$ ich nie ma.
\end{przyklad}

\begin{figure}[htbp]
\centering
\begin{tikzpicture}[>=Latex,font=\small]
  \node[font=\small] at (-1.9,1.25) {$\Z$};
  \node[font=\small] at (-1.9,-.05) {$\Z/3\Z$};
  \node[draw=manifoldblue!65!black,fill=manifoldblue!8,
        rounded corners=2pt,minimum width=3.1cm,minimum height=.7cm]
        (c0) at (0,1.25) {$\ldots,-3,0,3,6,\ldots$};
  \node[draw=manifoldblue!65!black,fill=manifoldblue!8,
        rounded corners=2pt,minimum width=3.1cm,minimum height=.7cm]
        (c1) at (3.5,1.25) {$\ldots,-2,1,4,7,\ldots$};
  \node[draw=manifoldblue!65!black,fill=manifoldblue!8,
        rounded corners=2pt,minimum width=3.1cm,minimum height=.7cm]
        (c2) at (7,1.25) {$\ldots,-1,2,5,8,\ldots$};
  \node[draw=manifoldblue!75!black,fill=manifoldblue!18,
        circle,minimum size=.8cm] (r0) at (0,-.05) {$[0]$};
  \node[draw=manifoldblue!75!black,fill=manifoldblue!18,
        circle,minimum size=.8cm] (r1) at (3.5,-.05) {$[1]$};
  \node[draw=manifoldblue!75!black,fill=manifoldblue!18,
        circle,minimum size=.8cm] (r2) at (7,-.05) {$[2]$};
  \draw[->,thick,manifoldblue!75!black] (c0)--(r0);
  \draw[->,thick,manifoldblue!75!black] (c1)--(r1);
  \draw[->,thick,manifoldblue!75!black] (c2)--(r2);
\end{tikzpicture}
\caption{Iloraz $\Z/3\Z$: wszystkie liczby w danym
górnym polu mają tę samą resztę, więc odwzorowanie ilorazowe
utożsamia je z jedną klasą.}
\label{fig:reszty-modulo-3-5}
\end{figure}

\section{Moduły, homomorfizmy i ilorazy}
\label{sec:moduly}

Przez $R$ oznaczamy pierścień przemienny z jedynką
$1\ne0$ w sensie Definicji~\ref{def:pierscien-przemienny}.
Pozwala to traktować wszystkie moduły
jednakowo, bez rozróżniania strony działania.
Wiele definicji działa również dla pierścieni
nieprzemiennych i modułów lewych; zaznaczymy,
gdzie przemienność upraszcza zapis.
Moduł można traktować jak przestrzeń liniową, w której
skalary nie muszą mieć odwrotności. W szczególności
wektory kratowe $\Z^2$ można dodawać i mnożyć przez
liczby całkowite, lecz połowa wektora $(1,0)$
nie należy do tej kraty.

\begin{definicja}[Moduł nad pierścieniem]
\index{modul nad pierscieniem@Moduł nad pierścieniem}
\label{def:modul}
\emph{Moduł $M$ nad $R$} to grupa abelowa
$(M,+,0)$ z mnożeniem przez skalary
$R\times M\to M$, $(r,m)\mapsto rm$, takim że
dla $r,s\in R$ i $m,n\in M$ mamy
\[
 (r+s)m=rm+sm,\qquad r(m+n)=rm+rn,\qquad
 (rs)m=r(sm),\qquad 1m=m.
\]
Grupa abelowa oznacza tutaj, że dodawanie w $M$
jest łączne, przemienne, ma element $0$ i każdy
element $m$ ma element przeciwny $-m$.
\end{definicja}

Aksjomaty wymuszają $0_Rm=0_M$ i $r0_M=0_M$:
na przykład $0_Rm=(0_R+0_R)m=0_Rm+0_Rm$,
więc skracamy w grupie addytywnej $M$; dla drugiej
równości stosujemy $r(0_M+0_M)=r0_M+r0_M$.
Ponadto $m+(-1_R)m=(1_R-1_R)m=0_M$,
czyli $(-1_R)m=-m$.

Definicja różni się od definicji przestrzeni liniowej
jedynie tym, że nie wymagamy dzielenia przez
niezerowy skalar. Ta różnica ma poważne skutki:
moduł może nie mieć bazy, a równanie $rm=0$ może
zachodzić dla $r\ne0$ i $m\ne0$.
W przypadku grup abelowych, czyli $\Z$-modułów,
opisujemy to precyzyjnie następująco.

\begin{definicja}[Element torsyjny i $m$-torsja]
\label{def:torsja}
Element $x$ grupy abelowej $M$ jest
\emph{torsyjny}, jeśli istnieje liczba całkowita
$k\ne0$ taka, że $kx=0$. Zbiór takich elementów
oznaczamy $\operatorname{Tors}(M)$ i nazywamy
\emph{podgrupą torsyjną}. Dla ustalonego $m\geq2$
definiujemy \emph{$m$-torsję}
$M[m]:=\{x\in M:mx=0\}=\ker(\times m:M\to M)$.
Podgrupa $M[m]$ składa się z elementów zabijanych
przez tę jedną liczbę $m$; cały
$\operatorname{Tors}(M)$ dopuszcza różne niezerowe
liczby dla różnych elementów.
\end{definicja}

Oba zbiory są rzeczywiście podgrupami: jeśli
$kx=0$ i $\ell y=0$, to $k\ell(x-y)=0$;
dla ustalonego $m$ mamy od razu
$m(x-y)=mx-my=0$. Na przykład
$\operatorname{Tors}(\Z\oplus\Z/2)
=0\oplus\Z/2$, podczas gdy $\Z$
nie ma niezerowej torsji.
W $M=\Z/6\Z$ wszystkie elementy są torsyjne,
ale $M[2]=\{[0],[3]\}$ i
$M[3]=\{[0],[2],[4]\}$. Na przykład
$2[k]=[0]$ dokładnie wtedy, gdy $6$ dzieli $2k$,
czyli $3$ dzieli $k$. Pokazuje to różnicę między
torsją o ustalonym wykładniku a całą podgrupą torsyjną.

\begin{przyklad}[Od grup abelowych do spinorów]
\label{prz:moduly-podstawowe}
\begin{enumerate}[leftmargin=*]
\item Każda przestrzeń liniowa nad ciałem
      $\mathbb K$ jest $\mathbb K$-modułem.
\item $R^n$ z dodawaniem i mnożeniem współrzędnych
      przez skalary jest wolnym $R$-modułem o bazie
      $e_1,\ldots,e_n$. W szczególności $\Z^2$
      jest modułem wolnym, choć nie jest przestrzenią
      liniową nad $\Z$, bo $\Z$ nie jest ciałem.
\item Każda grupa abelowa $A$ jest $\Z$-modułem:
      $na$ dla $n>0$ oznacza sumę $n$ kopii $a$,
      $0a=0$, a $(-n)a=-(na)$. Aksjomaty modułu
      wynikają z praw dodawania. Odwrotnie,
      w każdym $\Z$-module aksjomaty wymuszają
      właśnie tę regułę, więc pojęcia są równoważne.
\item Dla $m\geq2$ grupa $\Z/m\Z$ jest $\Z$-modułem;
      niezerowa klasa $[1]$ spełnia $m[1]=0$. Nie ma bazy
      nad $\Z$. Istotnie, gdyby $[1]=\sum_i a_i e_i$ w takiej
      bazie, równość $\sum_i ma_i e_i=0$ i jednoznaczność
      współczynników wymuszałyby $ma_i=0$ w $\Z$ dla każdego $i$.
      Stąd wszystkie $a_i=0$, sprzecznie z $[1]\ne0$.
\item $R/I$ dla ideału $I\subseteq R$ jest
      $R$-modułem: $r[s]=[rs]$. Działanie jest
      dobrze określone, gdyż z $s-s'\in I$
      wynika $r(s-s')\in I$.
      Sam ideał $I$ jest podmodułem modułu $R$.
\item Moduł spinorów z Definicji~\ref{def:spinor}
      jest modułem nad zespoloną algebrą
      Clifforda. Tam algebra może być
      nieprzemienna i działanie jest \emph{lewe};
      jest to przykład szerszego pojęcia modułu lewego.
      W głównej części obecnego rozdziału ograniczamy się
      do pierścieni przemiennych, wystarczających
      w podstawowych przykładach grup homologii.
\end{enumerate}
\end{przyklad}

\begin{definicja}[Podmoduł, homomorfizm i moduł wolny]
\index{podmodul homomorfizm modul wolny@Podmoduł, homomorfizm, moduł wolny}
\label{def:modul-map}
Podzbiór $N\subseteq M$ jest \emph{podmodułem}, jeśli zawiera $0$
i jest zamknięty na dodawanie, branie elementów przeciwnych
oraz mnożenie przez elementy $R$.
\emph{Homomorfizm $R$-modułów} $f:M\to N$ spełnia
$f(m+m')=f(m)+f(m')$ i $f(rm)=rf(m)$.
Zbiór takich homomorfizmów oznaczamy przez $\operatorname{Hom}_R(M,N)$.

Moduł \emph{wolny} z bazą $(e_i)_{i\in I}$ ma tę własność,
że każdy element jest jednoznaczną skończoną sumą
$\sum_{i\in I}r_i e_i$. Modelem jest $R^{(I)}$, czyli rodziny
współczynników o skończonym nośniku. Suma prosta $M\oplus N$
składa się z par $(m,n)$ z działaniami składowymi.
\end{definicja}

Przy dodawaniu punktowym i działaniu $(rf)(m)=r f(m)$
zbiór $\operatorname{Hom}_R(M,N)$ jest $R$-modułem.
Przemienność $R$ zapewnia $R$-liniowość $rf$:
$(rf)(sm)=rsf(m)=s(rf)(m)$.
Na przykład $\operatorname{Hom}_{\Z}(\Z/2,\Z)=0$,
bo obraz klasy $[1]$ musiałby być elementem $\Z$ zabijanym przez $2$.

\begin{stwierdzenie}[Homomorfizmy wyznaczone przez generator]
\label{prop:hom-cyclic-5}
Dla $R$-modułu $M$ odwzorowanie $f\mapsto f(1)$ daje
izomorfizm $\operatorname{Hom}_R(R,M)\cong M$.
Dla grupy abelowej $A$ i $n\geq2$ analogicznie
$\operatorname{Hom}_{\Z}(\Z/n\Z,A)\cong A[n]$,
przy czym izomorfizm wysyła $f$ na $f([1])$.
\end{stwierdzenie}
\begin{proof}
Liniowość daje $f(r)=rf(1)$, więc wartość w $1$
wyznacza $f:R\to M$. Odwrotnie, dla dowolnego $m\in M$
wzór $f_m(r)=rm$ definiuje homomorfizm z $f_m(1)=m$.
Oba przyporządkowania zachowują dodawanie i mnożenie
przez skalary, więc są odwrotnymi izomorfizmami.

Jeśli $f:\Z/n\Z\to A$, to
$nf([1])=f(n[1])=f([0])=0$, zatem $f([1])\in A[n]$.
Dla $x\in A[n]$ wzór $f_x([k])=kx$ jest dobrze określony:
gdy $[k]=[\ell]$, mamy $k-\ell=qn$ i
$kx-\ell x=qn x=0$. Wzór jest addytywny,
a $f_x([1])=x$. Znów dostajemy odwrotne homomorfizmy.
\end{proof}

Na przykład $\operatorname{Hom}_{\Z}(\Z/6\Z,\Z/4\Z)$
ma dwa elementy: obraz $[1]$ może być tylko $[0]$
albo $[2]$, bo tylko te klasy w $\Z/4\Z$ są
zabijane przez $6$. Istotnie, $6[k]_4=[0]_4$ oznacza
$4\mid6k$, równoważnie $2\mid k$; spośród reprezentantów
$0,1,2,3$ warunek spełniają dokładnie $0$ i $2$.

\begin{definicja}[Moduł ilorazowy]
\index{modul ilorazowy@Moduł ilorazowy}
\label{def:modul-ilorazowy}
Niech $M$ będzie $R$-modułem, a $N\subseteq M$ jego podmodułem.
Na $M$ wprowadzamy relację $m\sim m'$ wtedy i tylko wtedy,
gdy $m-m'\in N$. Jest to relacja równoważności:
zwrotność wynika z $0\in N$, symetria z zamknięcia na
elementy przeciwne, a przechodniość z zamknięcia na dodawanie.
Klasę elementu $m\in M$ oznaczamy przez $[m]=m+N$.
Zbiór $M/N:=\{m+N:m\in M\}$ nazywamy
\emph{modułem ilorazowym} $M$ przez $N$.
Klasa $m+N$ jest kopią podmodułu $N$ przesuniętą o $m$.
W ilorazie wszystkie jej punkty stają się jednym elementem.

Działania definiujemy przez $[m]+[m']=[m+m']$ oraz
$r[m]=[rm]$. Są dobrze określone. Istotnie, jeśli
$m-\widetilde m\in N$ i $m'-\widetilde m'\in N$, to
\[
(m+m')-(\widetilde m+\widetilde m')
=
(m-\widetilde m)+(m'-\widetilde m')
\in N,
\]
a dla $r\in R$ mamy $rm-r\widetilde m=r(m-\widetilde m)\in N$.
\end{definicja}

Naturalne odwzorowanie $\pi\colon M\to M/N$, $\pi(m)=[m]$,
jest surjektywnym homomorfizmem $R$-modułów o jądrze
$\ker\pi=N$.
Otrzymujemy więc krótki ciąg dokładny
\[
0\longrightarrow N
\overset{\iota}{\longrightarrow} M
\overset{\pi}{\longrightarrow}M/N
\longrightarrow0.
\]
Tutaj $\iota$ jest inkluzją. Słowo \emph{dokładny} oznacza,
że w każdym środkowym miejscu obraz poprzedniej strzałki
jest jądrem następnej: $\ker\iota=0$,
$\operatorname{im}\iota=N=\ker\pi$ oraz
$\operatorname{im}\pi=M/N$.

\begin{stwierdzenie}[Własność modułu ilorazowego]
\label{prop:module-quotient-universal-5}
Dla homomorfizmu $R$-modułów $f:M\to P$ istnieje
homomorfizm $\bar f:M/N\to P$ spełniający
$f=\bar f\circ\pi$ dokładnie wtedy, gdy
$N\subseteq\ker f$. Wtedy $\bar f$ jest jedyny
i $\bar f([m])=f(m)$.
\end{stwierdzenie}
\begin{proof}
Jeśli $f=\bar f\circ\pi$ i $n\in N$, to
$f(n)=\bar f([n])=\bar f([0])=0$.
Odwrotnie, jeśli $N\subseteq\ker f$, to z
$[m]=[m']$ wynika $m-m'\in N$ i
$f(m)-f(m')=f(m-m')=0$. Wzór na $\bar f$
jest więc dobrze określony. Zachowuje dodawanie
i mnożenie przez skalary, ponieważ robi to $f$.
Równość $f=\bar f\circ\pi$ zachodzi na każdym $m$,
a surjektywność $\pi$ zapewnia jednoznaczność $\bar f$.
\end{proof}

\begin{twierdzenie}[Pierwsze twierdzenie o izomorfizmie]
\label{thm:modul-izomorfizm}
Dla homomorfizmu $f:M\to P$ jądro i obraz są podmodułami,
a odwzorowanie
\[
\overline f:M/\ker f\longrightarrow\operatorname{im}f,
\qquad [m]\longmapsto f(m)
\]
jest izomorfizmem $R$-modułów.
\end{twierdzenie}
\begin{proof}
Jeśli $f(m)=f(m')=0$, to $f(m-m')=0$ i $f(rm)=0$;
zatem jądro jest podmodułem. Obraz jest nim, ponieważ
$f(m)+f(m')=f(m+m')$ oraz $rf(m)=f(rm)$.
Jeśli $[m]=[m']$ modulo $\ker f$, to $f(m-m')=0$,
więc $f(m)=f(m')$: wzór jest dobrze określony.
Jest liniowy i surjektywny na obraz. Równość
$\overline f([m])=0$ oznacza $m\in\ker f$, czyli $[m]=[0]$;
stąd jest również injektywny.
\end{proof}

\begin{przyklad}
Dla $M=\Z$ i $N=m\Z$ otrzymujemy $M/N=\Z/m\Z$.
Jest to moduł nad $\Z$, lecz dla $m>1$ nie jest modułem wolnym.
Przykład ten pokazuje, że iloraz modułu wolnego przez podmoduł
nie musi być modułem wolnym.

Dla $M=\R^2$ i $N=\R(1,0)$ odwzorowanie
$f(x,y)=y$ ma jądro $N$. Pierwsze twierdzenie
o izomorfizmie daje $\R^2/\R(1,0)\cong\R$;
izomorfizm wysyła $[(x,y)]$ na $y$.
Klasa punktu $(x,y)$ jest całą poziomą prostą
$\{(x+t,y):t\in\R\}$: w ilorazie zapamiętujemy tylko
wysokość $y$, jak na Rysunku~\ref{fig:modul-iloraz-poziome-5}.

Z kolei homomorfizm
$\Z^2\to\Z/2\Z\oplus\Z/3\Z$ dany przez
$(a,b)\mapsto([a]_2,[b]_3)$ jest surjektywny,
a jego jądrem jest $2\Z\oplus3\Z$. Stąd
$\Z^2/(2\Z\oplus3\Z)\cong\Z/2\Z\oplus\Z/3\Z$.
Ten iloraz jest skończony i torsyjny, mimo że
$\Z^2$ i jego podmoduł są wolne.
\end{przyklad}

\begin{figure}[htbp]
\centering
\begin{tikzpicture}[>=Latex,font=\small]
  \draw[->,black!70] (-1.9,0)--(1.9,0) node[right] {$x$};
  \draw[->,black!70] (0,-1.45)--(0,1.65) node[above] {$y$};
  \draw[manifoldblue!70!black,thick] (-1.55,-1)--(1.55,-1);
  \draw[manifoldblue!70!black,thick] (-1.55,0)--(1.55,0);
  \draw[manifoldblue!70!black,thick] (-1.55,1)--(1.55,1);
  \fill[manifoldblue!80!black] (-.95,1) circle (2pt);
  \fill[manifoldblue!80!black] (.8,1) circle (2pt);
  \node[above] at (-.95,1) {$(x_1,1)$};
  \node[above] at (.8,1) {$(x_2,1)$};
  \node[below] at (0,-1.45) {$\R^2$};
  \draw[->,thick,manifoldblue!75!black] (2.25,.2)--(3.55,.2)
        node[midway,above] {$[(x,y)]\mapsto y$};
  \draw[->,black!70] (4.55,-1.45)--(4.55,1.65) node[above] {$\R$};
  \foreach \v in {-1,0,1}{
    \fill[manifoldblue!80!black] (4.55,\v) circle (2.5pt);
  }
  \node[right] at (4.55,1) {$1$};
  \node[right] at (4.55,0) {$0$};
  \node[right] at (4.55,-1) {$-1$};
\end{tikzpicture}
\caption{Iloraz $\R^2/\R(1,0)$: każda pozioma prosta
jest jedną klasą. Dwa zaznaczone punkty na wysokości $1$
mają ten sam obraz.}
\label{fig:modul-iloraz-poziome-5}
\end{figure}

\begin{stwierdzenie}[Wymiar przestrzeni ilorazowej]
\index{modul ilorazowy!wymiar}
\label{prop:wymiar-ilorazu}
Niech $\mathbb K$ będzie ciałem, $M$ skończenie wymiarową
$\mathbb K$-przestrzenią wektorową, a $N\subseteq M$ podprzestrzenią.
Wówczas
\[
\dim_{\mathbb K}(M/N)
=
\dim_{\mathbb K}M-\dim_{\mathbb K}N.
\]
\end{stwierdzenie}

\begin{proof}
Niech $e_1,\ldots,e_k$ będzie bazą $N$.
Uzupełnijmy ją do bazy $e_1,\ldots,e_k,e_{k+1},\ldots,e_n$
przestrzeni $M$. Pokażemy, że klasy
$[e_{k+1}],\ldots,[e_n]$ tworzą bazę $M/N$.

Każdy $m\in M$ ma jednoznaczne współrzędne w tej bazie.
Pierwsza część jego rozkładu leży w $N$, więc
\[
\begin{aligned}
m&=\underbrace{\sum_{i=1}^k a_i e_i}_{\in N}
   +\sum_{i=k+1}^n a_i e_i,\\
[m]&=\sum_{i=k+1}^n a_i[e_i].
\end{aligned}
\]
Wskazane klasy generują zatem $M/N$.

Jeśli $\sum_{i=k+1}^n a_i[e_i]=[0]$, to
$v:=\sum_{i=k+1}^n a_i e_i\in N$.
Możemy więc zapisać $v=\sum_{i=1}^k b_i e_i$.
Jednoznaczność rozwinięcia względem pełnej bazy
wymusza $a_{k+1}=\cdots=a_n=0$.
Klasy są liniowo niezależne, więc tworzą bazę
$M/N$ o $n-k$ elementach. To daje żądany wzór.
\end{proof}

Dla modułów nad dziedziną całkowitości analogiczny wzór
sformułujemy po zdefiniowaniu iloczynu tensorowego modułów.

\section{Iloczyn tensorowy modułów}
\label{sec:tensor-modulow-algebra}
\index{iloczyn tensorowy modulow@Iloczyn tensorowy modułów}

W Rozdziale~\ref{ch:algebra-tensorowa} konstruowaliśmy iloczyny tensorowe
przestrzeni liniowych nad ciałem. Nad pierścieniem
budujemy podobny obiekt, lecz nie wolno już korzystać
z bazy dowolnego modułu. Intuicja pozostaje ta sama:
symbol $m\otimes n$ ma być liniowy w obu argumentach,
a każda dwuliniowa reguła na parach powinna przechodzić
przez jedno odwzorowanie liniowe.
Różnica względem przestrzeni liniowych ujawni się już w przykładzie:
utworzenie iloczynu tensorowego z $R/I$ narzuca w drugim czynniku relacje $im=0$
dla wszystkich $i\in I$, nawet jeśli wcześniej $im\ne0$.

\begin{definicja}[Iloczyn tensorowy modułów]
\label{def:tensor-modulow}
Niech $M,N$ będą $R$-modułami. Weźmy wolny
$R$-moduł $F=R^{(M\times N)}$ na formalnych
symbolach $[(m,n)]$, po jednym dla każdej pary $(m,n)$.
Nawias kwadratowy oznacza tu generator wolnego modułu;
dopiero $[(m,n)]+K$ będzie jego klasą w ilorazie.
Niech $K\subseteq F$ będzie
podmodułem generowanym przez wszystkie elementy
\[
\begin{aligned}
 &[(m+m',n)]-[(m,n)]-[(m',n)],\\
 &[(m,n+n')]-[(m,n)]-[(m,n')],\\
 &[(rm,n)]-r[(m,n)],\qquad [(m,rn)]-r[(m,n)],
\end{aligned}
\]
gdzie $m,m'\in M$, $n,n'\in N$ i $r\in R$.
\emph{Iloczyn tensorowy} to $M\otimes_R N=F/K$,
a $m\otimes n$ oznacza klasę $[(m,n)]+K$.
\end{definicja}

Podmoduł generowany przez relacje składa się
z ich \emph{skończonych} kombinacji $R$-liniowych.
To jest algebraiczny sens narzucenia relacji.
Ukończony iloczyn tensorowy jest innym pojęciem:
wymaga dodatkowej topologii lub normy na modułach
i nie występuje w tej konstrukcji.
Z definicji
$rm\otimes n=m\otimes rn=r(m\otimes n)$.
Ponadto $(0+0)\otimes n=0\otimes n+0\otimes n$,
więc $0\otimes n=0$; analogicznie $m\otimes0=0$.
Każdy element jest skończoną sumą tensorów
prostych, lecz nie musi być jednym tensorem prostym.
Przykład sumy niebędącej tensorem prostym pojawił się już
w Rozdziale~\ref{ch:algebra-tensorowa}; nad pierścieniem
nie możemy go na ogół rozpoznać za pomocą macierzy
współczynników, bo moduły nie muszą mieć baz.

\begin{twierdzenie}[Własność uniwersalna iloczynu tensorowego]
\label{thm:tensor-modulow-uniwersalnosc}
Dla każdego $R$-modułu $P$ i każdego
odwzorowania $R$-dwuliniowego
$b:M\times N\to P$ istnieje dokładnie
jeden homomorfizm $\widetilde b:M\otimes_R N\to P$
spełniający $\widetilde b(m\otimes n)=b(m,n)$.
\end{twierdzenie}
\begin{proof}
Na wolnym module $F$ istnieje dokładnie
jeden homomorfizm $L:F\to P$ z
$L([(m,n)])=b(m,n)$: określamy go na bazie
i przedłużamy liniowo. Addytywność i
$R$-liniowość $b$ w każdej zmiennej oznaczają,
że $L$ zeruje cztery rodzaje generatorów $K$.
Stąd $K\subseteq\ker L$, więc wzór
$\widetilde b(x+K)=L(x)$ jest dobrze określony.
Tensory proste rozpinają iloraz $F/K$;
wartości na nich wymuszają jednoznaczność.
Jeśli para $(T,\tau)$ również spełnia podaną własność,
stosujemy ją do odwzorowania $(m,n)\mapsto m\otimes n$
i otrzymujemy $\alpha:T\to M\otimes_R N$. Własność właśnie
skonstruowanego iloczynu daje w drugą stronę
$\beta:M\otimes_R N\to T$, wysyłające $m\otimes n$ na
$\tau(m,n)$. Złożenie $\beta\circ\alpha$ spełnia
$(\beta\circ\alpha)(\tau(m,n))=\tau(m,n)$, zatem
jest $\id_T$ przez jednoznaczność we własności $(T,\tau)$.
Analogicznie $\alpha\circ\beta$ jest $\id_{M\otimes_R N}$
przez właśnie dowiedzioną jednoznaczność.
To uzasadnia kanoniczność konstrukcji.
\end{proof}

\begin{przyklad}[Jednostka i iloraz przez ideał]
\label{prz:tensor-jednostka-ideal}
Odwzorowanie $R\otimes_R M\to M$,
$r\otimes m\mapsto rm$, jest izomorfizmem;
jego odwrotność wysyła $m$ na $1\otimes m$.
Oba złożenia sprawdzamy na generatorach:
$m\mapsto1\otimes m\mapsto m$ oraz
$r\otimes m\mapsto rm\mapsto1\otimes rm=r\otimes m$.
Ostatnia równość jest relacją przenoszącą skalar między czynnikami.
Jeśli $I\subseteq R$ jest ideałem, połóżmy
$IM=\{\sum_{j=1}^q i_jm_j:
i_j\in I,\ m_j\in M\}$.
Wówczas
\begin{equation}\label{eq:tensor-ideal-iloraz}
 (R/I)\otimes_R M\cong M/IM,
 \qquad (r+I)\otimes m\longmapsto rm+IM.
\end{equation}
Istotnie, wzór jest dwuliniowy i nie zależy
od wyboru $r$, bo $i m\in IM$ dla $i\in I$.
Odwrotność wysyła $m+IM$ na
$(1+I)\otimes m$. Jeśli $i\in I$, to
$(1+I)\otimes im=(i+I)\otimes m=0$;
znika więc na całym $IM$. Oba złożenia
są identycznościami na generatorach.
Dla $R=\Z$ i $I=m\Z$ otrzymujemy
$(\Z/m)\otimes_{\Z}A\cong A/mA$.
\end{przyklad}

\begin{przyklad}[Iloczyn tensorowy modułów cyklicznych]
Dla dodatnich liczb całkowitych $m,n$ wzór
\eqref{eq:tensor-ideal-iloraz} daje
\[
 (\Z/m)\otimes_{\Z}(\Z/n)
 \cong (\Z/n)/m(\Z/n)
 \cong \Z/(m\Z+n\Z)
 \cong \Z/\gcd(m,n).
\]
Środkowy izomorfizm pochodzi z surjekcji
$\Z/n\to\Z/(m\Z+n\Z)$, $[a]_n\mapsto a+(m\Z+n\Z)$.
Jest ona dobrze określona, bo $n\Z\subseteq m\Z+n\Z$.
Jej jądro to $m(\Z/n)$: warunek $a=mk+n\ell$
jest równoważny $[a]_n=m[k]_n$.
Stosujemy Twierdzenie~\ref{thm:modul-izomorfizm}.

Wyjaśnijmy też ostatnią równość ideałów. Na mocy
Stwierdzenia~\ref{prop:idealy-Z-5} niezerowy ideał
$m\Z+n\Z$ ma postać $d\Z$, gdzie $d>0$.
Ponieważ $m,n\in d\Z$, liczba $d$ dzieli oba.
Z kolei $d\in m\Z+n\Z$ daje $d=um+vn$
dla pewnych $u,v\in\Z$; każdy wspólny dzielnik $m,n$
dzieli więc $d$. Stąd $d=\gcd(m,n)$.
Jest to zarazem tożsamość Bézouta użyta w tym rachunku.
Na przykład $6(\Z/4)=\{[0]_4,[2]_4\}=2(\Z/4)$,
więc $(\Z/6)\otimes_{\Z}(\Z/4)\cong\Z/2$.
Natomiast $(\Z/2)\otimes_{\Z}(\Z/3)=0$:
mnożenie przez $2$ jest surjekcją na $\Z/3$,
gdyż klasy $[0]_3,[1]_3,[2]_3$ przechodzą odpowiednio
na $[0]_3,[2]_3,[1]_3$. Zatem $2(\Z/3)=\Z/3$.
\end{przyklad}

\begin{stwierdzenie}[Iloczyn tensorowy odwzorowań]
\label{prop:tensor-maps-algebra}
Homomorfizm $f:M\to M'$ indukuje
$f\otimes\id_N:M\otimes_R N\to M'\otimes_R N$
przez $(f\otimes\id_N)(m\otimes n)=f(m)\otimes n$.
Jeśli $f$ jest surjekcją, to $f\otimes\id_N$
jest surjekcją. Funktor $-\otimes_R N$ nie musi
zachowywać iniekcji.
\end{stwierdzenie}
\begin{proof}
Dwuliniowość odwzorowania $(m,n)\mapsto f(m)\otimes n$
i Twierdzenie~\ref{thm:tensor-modulow-uniwersalnosc}
dają indukowany homomorfizm.
Gdy $f$ jest surjekcją, każdy tensor prosty
$m'\otimes n$ ma przeciwobraz $m\otimes n$
z $f(m)=m'$. Tensory proste rozpinają cel,
więc indukowany homomorfizm jest surjekcją.
Dla iniekcji $f=\times2:\Z\to\Z$ oraz
$N=\Z/2$ identyfikacja
$\Z\otimes_{\Z}\Z/2\cong\Z/2$ pokazuje,
że $f\otimes\id_N$ jest odwzorowaniem zerowym:
$(2z)\otimes[1]=z\otimes[2]=0$.
Nie jest zatem iniekcją.
\end{proof}

\begin{stwierdzenie}[Dokładność po prawej stronie]
\label{prop:tensor-prawostronnie-dokladny}
\index{iloczyn tensorowy!dokładność po prawej stronie}
Jeśli ciąg $R$-modułów
\[
 M_1\xrightarrow{f}M_2\xrightarrow{q}M_3\longrightarrow0
\]
jest dokładny, to dla każdego $R$-modułu $N$ dokładny jest także
\[
 M_1\otimes_R N\xrightarrow{f\otimes\id_N}
 M_2\otimes_R N\xrightarrow{q\otimes\id_N}
 M_3\otimes_R N\longrightarrow0.
\]
Nie ma analogicznej gwarancji iniektywności na lewym końcu.
\end{stwierdzenie}
\begin{proof}
Ponieważ $q$ jest surjekcją, taka jest również $q\otimes\id_N$.
Równość $q\circ f=0$ daje
$\operatorname{im}(f\otimes\id_N)\subseteq\ker(q\otimes\id_N)$.
Oznaczmy iloraz $(M_2\otimes_R N)/\operatorname{im}(f\otimes\id_N)$
przez $Q$. Dla $x\in M_3$ wybierzmy podniesienie
$\widehat x\in M_2$ i połóżmy
$b(x,n)=[\widehat x\otimes n]\in Q$.
Jeśli $\widehat x'$ jest innym podniesieniem, to
$\widehat x-\widehat x'\in\ker q=\operatorname{im}f$;
stąd obie klasy są równe. Podniesienia sumy oraz
sumy podniesień różnią się elementem $\ker q$,
podobnie jak podniesienia $rx$ oraz $r\widehat x$.
Dlatego $b$ jest $R$-dwuliniowe. Własność uniwersalna
daje homomorfizm $V:M_3\otimes_R N\to Q$.
W drugą stronę $q\otimes\id_N$ zeruje
$\operatorname{im}(f\otimes\id_N)$ i schodzi do
$U:Q\to M_3\otimes_R N$.
Dla $x\otimes n$ złożenie daje
$(U\circ V)(x\otimes n)=q(\widehat x)\otimes n=x\otimes n$.
Z kolei $(V\circ U)([m\otimes n])=[m\otimes n]$,
bo jako podniesienie $q(m)$ możemy wybrać samo $m$.
Takie elementy rozpinają odpowiednie moduły,
więc $U,V$ są wzajemnie odwrotne.
W szczególności $q\otimes\id_N$ jest złożeniem projekcji
na $Q$ z izomorfizmem $U$, a jego jądro jest jądrem tej projekcji.
To dokładnie równość
$\ker(q\otimes\id_N)=\operatorname{im}(f\otimes\id_N)$.
\end{proof}

\begin{definicja}[Rozszerzenie skalarów]
\label{def:rozszerzenie-skalarow}
\index{rozszerzenie skalarów}
Jeśli $\varphi:R\to S$ jest homomorfizmem
pierścieni przemiennych z jedynką, to $S$
staje się $R$-modułem przez $r\cdot s=\varphi(r)s$.
Na $M\otimes_R S$ działa $S$ przez
$s\cdot(m\otimes a)=m\otimes sa$.
Ten $S$-moduł nazywamy \emph{rozszerzeniem
skalarów} modułu $M$ wzdłuż $\varphi$.
\end{definicja}

Sprawdźmy relację zbalansowania dla działania
$S$: dla $r\in R$ oraz $s,a\in S$ mamy
\[
 s\cdot((rm)\otimes a)
 =m\otimes\varphi(r)sa
 =m\otimes s\varphi(r)a
 =s\cdot(m\otimes\varphi(r)a).
\]
Przemienność $S$ uzasadnia środkową równość;
relacje addytywności sprawdza się składnikowo.
Dlatego działanie nie zależy od zapisu elementu
jako sumy tensorów prostych. Dla wolnego
$M=\bigoplus_{i\in J}Re_i$ mamy
\[
 M\otimes_R S\cong\bigoplus_{i\in J}Se_i,
 \quad
 (\textstyle\sum_i r_i e_i)\otimes s
 \longmapsto\textstyle\sum_i\varphi(r_i)s e_i.
\]
Odwrotność wysyła $\sum_i s_i e_i$ na
$\sum_i e_i\otimes s_i$; wszystkie sumy są
skończone. Ogólniej, dla rodziny $R$-modułów $(M_i)_{i\in J}$
otrzymujemy izomorfizm
\[
 \left(\bigoplus_{i\in J}M_i\right)\otimes_R N
 \longrightarrow \bigoplus_{i\in J}(M_i\otimes_R N),
 \qquad (m_i)_i\otimes n\longmapsto(m_i\otimes n)_i.
\]
Wzór jest dwuliniowy i ma skończony nośnik, więc określa
homomorfizm przez własność uniwersalną. Odwrotność na
$i$-tym składniku wysyła $m\otimes n$ na $\iota_i(m)\otimes n$,
gdzie $\iota_i$ jest inkluzją do sumy prostej. Oba złożenia są
identycznościami na tensorach prostych, które rozpinają moduły.
Funktor $-\otimes_R N$ zachowuje też izomorfizmy: dla izomorfizmu $f$
odwrotnością $f\otimes\id$ jest $f^{-1}\otimes\id$.
Przykład z $\times2$ ostrzega natomiast przed
wnioskowaniem o dowolnych iniekcjach.
Na przykład po rozszerzeniu skalarów przez
$\R\hookrightarrow\C$ moduł $\R^2$ staje się
$\C^2$: nowe współczynniki są zespolone,
a dwa wektory bazowe pozostają niezależne.

\begin{uwaga}[Moduły rzutowe]
\label{uwaga:moduly-projektywne}
\index{moduł rzutowy}
Moduł $P$ jest \emph{rzutowy}, gdy dla każdej
surjekcji $q:E\twoheadrightarrow B$ i każdego
homomorfizmu $g:P\to B$ istnieje podniesienie
$\widetilde g:P\to E$ z $q\circ\widetilde g=g$.
Równoważnie, $P$ jest składnikiem prostym
pewnego modułu wolnego. Jeśli $P$ jest rzutowy,
wybieramy wolny moduł $F$ na symbolach elementów
$P$ i jego naturalną surjekcję $q:F\twoheadrightarrow P$.
Podnosimy $\id_P$ przez $q$, otrzymując $j:P\to F$
z $q\circ j=\id_P$. Izomorfizm i jego odwrotność to
\[
 P\oplus\ker q\longrightarrow F,\quad (p,k)\longmapsto j(p)+k,
 \qquad
 f\longmapsto\bigl(q(f),f-j(q(f))\bigr).
\]
Drugi składnik ostatniej pary należy do $\ker q$, bo
$q(f-j(q(f)))=q(f)-q(f)=0$.
Złożenie w kierunku $F\to F$ daje $j(q(f))+f-j(q(f))=f$,
a w kierunku $P\oplus\ker q$ daje $(p,k)$, bo $q(k)=0$.
To wyprowadza rozkład $F\cong P\oplus\ker q$. Odwrotnie,
odwzorowanie z wolnego modułu można podnieść,
wybierając przeciwobrazy wartości na bazie.
Jeśli $F=P\oplus Q$, to dla $g:P\to B$
podnosimy $g\circ\operatorname{pr}_P:F\to B$
i ograniczamy otrzymane odwzorowanie do $P$.
Moduł rzutowy nie musi być wolny: dla
$R=\Z/6$ ideał $3R=\{0,3\}$ jest składnikiem
prostym w $R=3R\oplus4R$, lecz jego dwa
elementy wykluczają izomorfizm z $R^{(J)}$.
Rzeczywiście, $3R\cap4R=\{0\}$, a $1=3+4$ w $R$ daje
$r=3r+4r$ dla każdego $r\in R$. Niezerowy moduł wolny
nad $R$ zawiera co najmniej sześć elementów, bo zawiera kopię $R$.
\end{uwaga}

W zastosowaniach do nieprzemiennego pierścienia
grupowego $\Z[\pi]$ rozróżnia się moduł prawy
i lewy. Dla dowolnego pierścienia z jedynką $\Lambda$
oraz prawego modułu $M$ i lewego $N$ zaczynamy od
\emph{wolnej grupy abelowej} na parach $(m,n)$.
Dzielimy ją przez podgrupę generowaną przez relacje
addytywności w obu argumentach oraz
$[(m\lambda,n)]-[(m,\lambda n)]$, $\lambda\in\Lambda$.
Otrzymana grupa abelowa $M\otimes_\Lambda N$ spełnia
$(m\lambda)\otimes n=m\otimes(\lambda n)$.
Każde odwzorowanie do grupy abelowej, addytywne w obu
argumentach i spełniające tę relację, przechodzi przez
dokładnie jeden homomorfizm z iloczynu tensorowego: przedłużamy je
na wolną grupę i dzielimy przez relacje, jak wyżej.
Bez dodatkowych danych wynik nie ma naturalnej struktury
$\Lambda$-modułu. Dla przemiennego $R$ utożsamiamy działanie
prawe z lewym przez $mr=rm$ i odzyskujemy wcześniejszy
$R$-moduł z działaniem $r(m\otimes n)=(rm)\otimes n$.
Na przykład dla prawego modułu regularnego $\Lambda$ zachodzi
$\Lambda\otimes_\Lambda N\cong N$:
$\lambda\otimes n\mapsto\lambda n$, a odwrotność
wysyła $n$ na $1\otimes n$. Relacja zbalansowania daje
$\lambda\otimes n=1\otimes\lambda n$, więc obydwa
złożenia są identycznościami.
Charakteryzacja rzutowości przez podnoszenie i
składniki proste wolnych działa również tam.

\section{Ranga modułu i ciało ułamków}
\label{sec:ranga-modulu}

Po rozszerzeniu skalarów z dziedziny całkowitości $R$ do jej
ciała ułamków moduł staje się przestrzenią wektorową. Jej wymiar
mierzy liczbę niezależnych kierunków, które pozostają po usunięciu
torsji.

\subsection*{Konstrukcja ciała ułamków}

\begin{definicja}[Ciało ułamków]
\label{def:cialo-ulamkow-5}
\index{ciało ułamków}
Niech $R$ będzie dziedziną całkowitości i $S=R\setminus\{0\}$.
Chcemy umożliwić dzielenie przez każdy element $S$.
Na parach $(a,s)\in R\times S$ określamy relację
$(a,s)\sim(b,t)$ wtedy i tylko wtedy, gdy $ta=sb$.
Klasę oznaczamy $a/s$, a zbiór klas przez $K=\operatorname{Frac}(R)$.
Dodawanie i mnożenie mają znane z liczb wymiernych postacie
\[
 \frac as+\frac bt=\frac{ta+sb}{st},\qquad
 \frac as\frac bt=\frac{ab}{st}.
\]
\end{definicja}

\begin{proof}[Poprawność konstrukcji]
Zbiór $S$ zawiera $1$ i jest zamknięty na mnożenie, ponieważ $R$
nie ma dzielników zera. Mianowniki pozostają więc dozwolone.
Zwrotność i symetria relacji są bezpośrednie. Jeśli $ta=sb$ oraz
$ub=tc$, to $t(ua-sc)=0$; skrócenie $t\ne0$ daje $ua=sc$.
Stąd przechodniość.

Jeśli $a/s=a'/s'$ i $b/t=b'/t'$, mamy $s'a=sa'$ oraz $t'b=tb'$.
Po sprowadzeniu sum do wspólnego mianownika korzystamy
w pierwszym składniku z $s'a=sa'$, a w drugim z $t'b=tb'$:
\[
 \begin{aligned}
 s't'(ta+sb)
 &=tt'(s'a)+ss'(t'b)\\
 &=stt'a'+sts'b'=st(t'a'+s'b'),\\
 s't'ab&=(s'a)(t'b)=(sa')(tb')=sta'b'.
 \end{aligned}
\]
Druga równość sprawdza zarazem niezależność iloczynu od przedstawicieli.
Prawa pierścienia wynikają z praw w $R$ po sprowadzeniu do wspólnego
mianownika. Na przykład suma trzech ułamków ma licznik
$tu\,a+su\,b+st\,c$ i mianownik $stu$ niezależnie od nawiasów.
Rozdzielność wynika z rachunku
\[
 \frac as\left(\frac bt+\frac cu\right)
 =\frac{a(ub+tc)}{stu}
 =\frac{uab+tac}{stu}
 =\frac{ab}{st}+\frac{ac}{su}.
\]
W ostatnim kroku obie sumowane klasy sprowadzono do mianownika $stu$.
Zerem jest $0/1$, jedynką $1/1$, a elementem przeciwnym do $a/s$ jest $(-a)/s$.

Równość $a/s=0/1$ zachodzi dokładnie dla $a=0$.
Dla niezerowej klasy $a/s$ element $s/a$ jest więc dozwolony
i $(a/s)(s/a)=1$. Ponadto $0/1\ne1/1$. Otrzymaliśmy ciało.
Odwzorowanie $R\to K$, $a\mapsto a/1$, zachowuje działania i jest
iniektywne; odtąd utożsamiamy $R$ z jego obrazem.
\end{proof}

\begin{stwierdzenie}[Własność uniwersalna ciała ułamków]
Jeśli $j:R\hookrightarrow L$ jest iniektywnym homomorfizmem
do ciała $L$, to istnieje dokładnie jeden homomorfizm ciał
$\widehat j:K\to L$ przedłużający $j$. Jest on dany wzorem
$\widehat j(a/s)=j(a)j(s)^{-1}$.
Jego obraz jest najmniejszym podciałem $L$ zawierającym $j(R)$.
\end{stwierdzenie}
\begin{proof}
Ponieważ $s\ne0$ i $j$ jest iniektywne, $j(s)$ ma odwrotność.
Jeśli $a/s=b/t$, to $ta=sb$, a więc
$j(t)j(a)=j(s)j(b)$; po podzieleniu przez $j(s)j(t)$
otrzymujemy równość wartości wzoru. Dodawanie i mnożenie
ułamków pokazują, że $\widehat j$ jest homomorfizmem.
Każdy homomorfizm ciał przedłużający $j$ musi wysłać
$a/s=(a/1)(s/1)^{-1}$ na $j(a)j(s)^{-1}$,
więc jest jedyny. Obraz składa się z takich ilorazów;
każde podciało zawierające $j(R)$ musi je zawierać.
\end{proof}

\begin{przyklad}[Liczby i funkcje wymierne]
Dla $R=\Z$ konstrukcja daje $\operatorname{Frac}(\Z)=\Q$.
Dla $R=\mathbb K[t]$ otrzymujemy ciało funkcji wymiernych
$\mathbb K(t)$: elementy to ilorazy wielomianów $p(t)/q(t)$ z
$q\ne0$, przy równości przez mnożenie na krzyż. Nie wymagamy,
by mianownik nie miał zer; chodzi o ułamki formalne, nie o funkcje
określone w każdym punkcie $\mathbb K$.
\end{przyklad}

\subsection*{Lokalizacja modułu i jej związek z iloczynem tensorowym}
\index{lokalizacja modułu}

Niech $S=R\setminus\{0\}$. Aby opisać elementy iloczynu tensorowego przez
ułamki, definiujemy \emph{lokalizację} $S^{-1}M$ jako zbiór klas
par $(m,s)\in M\times S$ przy relacji
\[
(m,s)\sim(m',s')
\quad\Longleftrightarrow\quad
t(s'm-sm')=0\text{ dla pewnego }t\in S.
\]
Klasę pary $(m,s)$ oznaczamy $m/s$. Mianowniki stają się tu
odwracalne: dla $M=R$ otrzymujemy $S^{-1}R=K$.
Relacja jest zwrotna i symetryczna. Jeśli
$u(s'm-sm')=0$ i $v(s''m'-s'm'')=0$, to
$uvs'(s''m-sm'')=0$: po rozwinięciu lewa strona
jest sumą $uvs''(s'm-sm')$ oraz
$uvs(s''m'-s'm'')$. Ponieważ $uvs'\in S$,
otrzymujemy przechodniość.
Dodawanie i działanie $K$ określamy przez
\[
\frac{m}{s}+\frac{m'}{s'}=\frac{s'm+sm'}{ss'},
\qquad
\frac{a}{u}\cdot\frac{m}{s}=\frac{am}{us}.
\]
Wzory nie zależą od przedstawicieli. Dla dodawania wynika to
z tożsamości
\[
tt'(s'm+sm')-ss'(t'n+tn')
=t's'(tm-sn)+ts(t'm'-s'n'),
\]
gdy $m/s=n/t$ oraz $m'/s'=n'/t'$; oba składniki po prawej
stronie zeruje iloczyn odpowiednich elementów $S$.
Dokładniej, jeśli $u(tm-sn)=0$ i $v(t'm'-s'n')=0$,
to $uv$ zeruje całą różnicę. Nie trzeba skracać elementów
$S$ w module $M$.
Jeśli ponadto $a/u=b/v$ w $K$, to $va=ub$ i różnica
krzyżowa dla $am/(us)$ i $bn/(vt)$ wynosi
$v a t m-u b s n=v a(t m-s n)$; zeruje ją ten sam element $S$
co różnicę $t m-s n$. Obejmuje to jednoczesną zmianę
przedstawicieli skalara i wektora, więc działanie $K$
jest dobrze określone.

W tym działaniu $a/u$ jest \emph{skalarem} z $K$, a $m/s$ jest
\emph{wektorem} z $S^{-1}M$; nie mnożymy dwóch elementów modułu.
Na $M\otimes_R K$ ten sam skalar działa wzorem
$(a/u)\cdot(m\otimes b/s)=m\otimes ab/(us)$.
Łączność działania na lokalizacji widać z równości
\[
 (a/u)((b/v)(m/s))=(abm)/(uvs)=((a/u)(b/v))(m/s).
\]
Dodawanie jest łączne: obydwa sposoby dodania trzech klas
$m/s$, $m'/s'$ i $m''/s''$ mają przedstawiciela
$(s's''m+ss''m'+ss'm'')/(ss's'')$.
Zerem jest $0/1$, a przeciwną klasą do $m/s$ jest $(-m)/s$.
Na przykład rozdzielność względem dodawania skalarów ma postać
\[
 \left(\frac au+\frac bv\right)\frac ms
 =\frac{(va+ub)m}{uvs}
 =\frac{am}{us}+\frac{bm}{vs}.
\]
Rozdzielność względem dodawania wektorów sprawdza się tak samo
po sprowadzeniu do wspólnego mianownika; $1/1$ działa
identycznie. Dostajemy więc przestrzeń wektorową nad $K$.
W przypadku $M=R$ relacja lokalizacji redukuje się do mnożenia
na krzyż, bo niezerowy $t$ można skrócić w $R$.
W ogólnym module tego robić nie wolno: w $\Z$-module $\Z/2\Z$
element $[1]$ jest niezerowy, lecz $[1]/1=0$ w lokalizacji,
gdyż $2[1]=0$. Właśnie dlatego relacja zawiera dodatkowy czynnik $t$.

\begin{przyklad}[Dwie lokalizacje nad liczbami całkowitymi]
Dla $M=\Z$ mamy $S^{-1}M=\Q$ i na przykład
$(2/3)\cdot(5/7)=10/21$.
Dla $M=\Z\oplus\Z/2\Z$ torsyjny składnik znika:
odwzorowanie $S^{-1}M\to\Q$ określone przez
$(a,[b])/s\mapsto a/s$
jest izomorfizmem, z odwrotnością $a/s\mapsto(a,[0])/s$.
Istotnie, różnicę $(0,[b])/s$ zeruje $2$, a równość dwóch
pierwszych współrzędnych sprawdza się przez mnożenie na krzyż.
\end{przyklad}

\begin{stwierdzenie}[Iloczyn tensorowy z ciałem ułamków]
\label{prop:tensor-cialo-ulamkow}
Niech $R$ będzie dziedziną całkowitości,
$S=R\setminus\{0\}$ i $K=\operatorname{Frac}(R)$.
Dla każdego $R$-modułu $M$ istnieje naturalny izomorfizm
\[
M\otimes_R K\cong S^{-1}M,
\qquad m\otimes\frac{a}{s}\longmapsto\frac{am}{s}.
\]
W szczególności $m/s=0$ dokładnie wtedy, gdy
$tm=0$ dla pewnego $t\in S$.
\end{stwierdzenie}
\begin{proof}
Z definicji lokalizacji
$m/s=0/1$ dokładnie wtedy, gdy pewien $t\in S$
zeruje $m$. Podany wzór jest $R$-dwuliniowy i nie zależy
od wyboru przedstawiciela $a/s\in K$: z $ta=sb$ wynika
$t(am)=s(bm)$. Dlatego przez
Twierdzenie~\ref{thm:tensor-modulow-uniwersalnosc}
określa homomorfizm $M\otimes_R K\to S^{-1}M$.
W drugą stronę kładziemy $m/s\mapsto m\otimes(1/s)$.
Jeśli $(m,s)\sim(m',s')$, to dla pewnego $t\in S$
zachodzi $t(s'm-sm')=0$, a więc
\[
m\otimes\frac1s-m'\otimes\frac1{s'}
=(s'm-sm')\otimes\frac1{ss'}
=t(s'm-sm')\otimes\frac1{tss'}=0.
\]
Odwzorowanie odwrotne jest więc dobrze określone.
Wzór na dodawanie ułamków pokazuje jego addytywność;
zgodność z działaniem $K$ widać na przykład z rachunku
\[
 \frac{am}{us}\longmapsto (am)\otimes\frac1{us}
 =m\otimes\frac{a}{us}
 =\frac au\cdot\left(m\otimes\frac1s\right).
\]
Oba złożenia sprawdzamy bezpośrednio:
$m/s\mapsto m\otimes(1/s)\mapsto m/s$ oraz
$m\otimes(a/s)\mapsto(am)/s\mapsto(am)\otimes(1/s)=m\otimes(a/s)$.
W drugiej równości użyto relacji przenoszącej skalar $a$.
Tensory proste generują cały iloczyn tensorowy,
więc złożenia są identycznościami.
Ponadto dla homomorfizmu $f:M\to M'$ oba sposoby
przeniesienia klasy $m/s$ dają $f(m)/s$;
izomorfizm jest więc naturalny.
\end{proof}

\subsection*{Ranga i jej addytywność}

\begin{definicja}[Ranga modułu nad dziedziną całkowitości]
\index{ranga modułu}
\label{def:ranga-modulu}
Niech $R$ będzie dziedziną całkowitości, a
$K=\operatorname{Frac}(R)$ jej ciałem ułamków.
Rangą $R$-modułu $M$ nazywamy liczbę
\[
\operatorname{rank}_R M
:=
\dim_K(M\otimes_R K),
\]
o ile wymiar po prawej stronie jest skończony.
\end{definicja}

Jeśli $M$ jest wolnym modułem o $n$ generatorach bazowych, to
$M\cong R^n$ i wzór na rozszerzenie skalarów daje
$M\otimes_R K\cong K^n$. Zatem $\operatorname{rank}_R M=n$.

\begin{przyklad}[Ranga nie widzi torsji]
Dla $M=\Z^2\oplus\Z/6$ lokalizacja daje
\[
 M\otimes_{\Z}\Q\longrightarrow\Q^2,\qquad
 ((a,b),[c])\otimes q\longmapsto(qa,qb).
\]
To izomorfizm: jego odwrotność wysyła $(x,y)$ na
$((1,0),[0])\otimes x+((0,1),[0])\otimes y$,
a dla $w=((0,0),[c])$ mamy
$w\otimes q=(6w)\otimes(q/6)=0$, bo $6[c]=0$.
Stąd $\operatorname{rank}_{\Z}M=2$.
Podobnie moduł $\mathbb K[t]/(t)$ ma rangę $0$
nad $\mathbb K[t]$: jego generator $[1]$ jest niezerowy,
ale $t[1]=0$, a $t$ staje się odwracalny w
$\mathbb K(t)$.
\end{przyklad}

\begin{stwierdzenie}[Addytywność rangi]
\index{ranga modułu!addytywnosc@addytywność}
\label{prop:ranga-ilorazu}
Niech $R$ będzie dziedziną całkowitości oraz
$N\subseteq M$ podmodułem. Jeśli odpowiednie rangi są skończone, to
\[
\operatorname{rank}_R M
=
\operatorname{rank}_R N
+
\operatorname{rank}_R(M/N).
\]
W szczególności $\operatorname{rank}_R(M/N)
=\operatorname{rank}_R M-\operatorname{rank}_R N$.
\end{stwierdzenie}

\begin{proof}
Z krótkiego ciągu dokładnego
\[
0\longrightarrow N
\longrightarrow M
\longrightarrow M/N
\longrightarrow0
\]
po utworzeniu iloczynu tensorowego każdego składnika
z ciałem ułamków $K=\operatorname{Frac}(R)$
otrzymujemy ciąg dokładny przestrzeni wektorowych
\[
0
\longrightarrow
N\otimes_R K
\longrightarrow
M\otimes_R K
\longrightarrow
(M/N)\otimes_R K
\longrightarrow0.
\]
Wyjaśnijmy tu krok o dokładności, korzystając z
izomorfizmu~\ref{prop:tensor-cialo-ulamkow}. Jeśli $n/s$
dla $n\in N$ znika w $M\otimes_R K$, to $tn=0$
dla pewnego $t\in S$ już w $N$, więc $n/s=0$ także
w $N\otimes_R K$. Jeśli obraz $m/s$ znika w
$(M/N)\otimes_R K$, to dla pewnego $t\in S$ mamy
$tm\in N$, a więc $m/s=(tm)/(ts)$ należy do obrazu
$N\otimes_R K$. Surjektywność na końcu wynika z
podnoszenia licznika klasy $[m]/s$.
Z Twierdzenia~\ref{thm:modul-izomorfizm} ostatnia przestrzeń
jest izomorficzna z ilorazem $M\otimes_R K$
przez obraz $N\otimes_R K$. Iniektywność pierwszej strzałki
utożsamia ten obraz z $N\otimes_R K$.
Stosując wzór na wymiar ilorazu ze
Stwierdzenia~\ref{prop:wymiar-ilorazu} i przenosząc składnik, dostajemy
\[
\dim_K(M\otimes_R K)
=
\dim_K(N\otimes_R K)
+
\dim_K((M/N)\otimes_R K),
\]
co jest dokładnie żądanym wzorem.
\end{proof}

\begin{przyklad}[Torsja powstała w ilorazie]
Weźmy $M=\Z^2$ oraz podmoduł
$N=2\Z\oplus\{0\}$, którego ranga wynosi $1$.
Odwzorowanie
\[
 M/N\longrightarrow\Z/2\oplus\Z,\qquad
 (a,b)+N\longmapsto([a]_2,b)
\]
jest izomorfizmem: dwie pary mają ten sam obraz
dokładnie wtedy, gdy ich różnica należy do $N$.
Iloraz ma więc niezerowy składnik torsyjny, ale jego
ranga wynosi $1$, zgodnie z równością $2=1+1$.
\end{przyklad}

\begin{uwaga}
Dla ogólnego pierścienia $R$ nie należy bez dodatkowych założeń
pisać wzoru $\dim(M/N)=\dim M-\dim N$.
Pojęcie wymiaru jest naturalne dla przestrzeni wektorowych, natomiast
dla modułów stosuje się między innymi rangę, liczbę generatorów
oraz -- w przypadku skończenie generowanych modułów nad
pierścieniami głównymi -- rozkład
na część wolną i torsyjną.

Na przykład $\Z/2\Z$ ma rangę zerową jako $\Z$-moduł,
mimo że jest modułem niezerowym. Ranga nie wykrywa więc torsji.
\end{uwaga}

\par\smallskip
{\footnotesize\noindent\emph{Dalsza lektura:}
Konstrukcję iloczynu wolnego i jego związek z topologią opisuje
Allen Hatcher, \emph{Algebraic Topology}, sekcja~1.2,
\emph{Free Products of Groups}, s.~41--42,
\url{https://pi.math.cornell.edu/~hatcher/AT/AT.pdf}.
Iloczyn tensorowy nad pierścieniem przemiennym rozwija
\emph{The Stacks Project}, sekcja~10.12, \emph{Tensor products},
\url{https://stacks.math.columbia.edu/tag/00CV}.
Tam też znajdują się dokładność po prawej stronie
(Lemat~10.12.10) i związek z lokalizacją (Lemat~10.12.15).
Lokalizację opisuje sekcja~10.9 tego samego źródła,
\url{https://stacks.math.columbia.edu/tag/00CM}.
Jej dokładność jest sformułowana w Propozycji~10.9.12,
\url{https://stacks.math.columbia.edu/tag/00CS}.
Powyższe konstrukcje i dowody nie wymagają jeszcze algebry homologicznej.\par}
